
\Data\ID{1003020001}
\Updated{2020-07-15 03:07:12}
\Level{A}
\SubArea{Quadratic equations and inequalities}
\LANGen{
\Question{
Find the sum of the roots of the equation.
\[
4x^2+2x=0
\]}
{\(-\frac12\)}{\(0\)}{\(2\)}{\(-2\)}{}{}}
\LANGpl{
\Question{
Oblicz sumę pierwiastków podanego równania.
\[
4x^2+2x=0
\]}
{\(-\frac12\)}{\(0\)}{\(2\)}{\(-2\)}{}{}}
\LANGcs{
\Question{
Určete součet kořenů rovnice.
\[
4x^2+2x=0
\]}
{\(-\frac12\)}{\(0\)}{\(2\)}{\(-2\)}{}{}}
\LANGsk{
\Question{
Určte súčet koreňov rovnice.
\[
4x^2+2x=0
\]}
{\(-\frac12\)}{\(0\)}{\(2\)}{\(-2\)}{}{}}
\LANGes{
\Question{
Determina la suma de las raíces de la ecuación.
\[
4x^2+2x=0
\]}
{\(-\frac12\)}{\(0\)}{\(2\)}{\(-2\)}{}{}}
\End

\Data\ID{1003047001}
\Updated{2020-08-20 07:08:36}
\Level{A}
\SubArea{Quadratic equations and inequalities}
\LANGen{
\Question{
We are given the equation \( 2x^2+10x=8x+2x^2 \).  Decide which of the following equations has the different set of roots than the given equation has, i.e. choose the equation which is not equivalent to the given equation.}
{\( 2x+10=8+2x \)}{\( 10x=8x \)}{\( 2x^2+2x=2x^2 \)}{\( x^2+5x=4x+x^2 \)}{}{}}
\LANGpl{
\Question{
Dane jest równanie \( 2x^2+10x=8x+2x^2 \).  Zdecyduj, które z poniższych równań posiada inny zbiór pierwiastków niż podane równanie, tzn wybierz równanie, które nie jest równoważne  podanemu równaniu.}
{\( 2x+10=8+2x \)}{\( 10x=8x \)}{\( 2x^2+2x=2x^2 \)}{\( x^2+5x=4x+x^2 \)}{}{}}
\LANGcs{
\Question{
Je dána rovnice	   \( 2x^2+10x=8x+2x^2 \).  Z následujících rovnic vyberte tu, která má jinou množinu kořenů, než zadaná rovnice, tj. není s danou rovnicí ekvivalentní.}
{\( 2x+10=8+2x \)}{\( 10x=8x \)}{\( 2x^2+2x=2x^2 \)}{\( x^2+5x=4x+x^2 \)}{}{}}
\LANGsk{
\Question{
Je daná rovnica \( 2x^2+10x=8x+2x^2 \).  Z nasledujúcich rovníc vyberte tú, ktorá má inú množinu koreňov, než zadaná rovnica, t.j. nie je s danou rovnicou ekvivalentná.}
{\( 2x+10=8+2x \)}{\( 10x=8x \)}{\( 2x^2+2x=2x^2 \)}{\( x^2+5x=4x+x^2 \)}{}{}}
\LANGes{
\Question{
Dada la ecuación \( 2x^2+10x=8x+2x^2 \).  Decide cuál de las siguientes ecuaciones tiene conjunto de soluciones diferente a la ecuación dada, o sea, elige la ecuación que no es equivalente a la ecuación dada.}
{\( 2x+10=8+2x \)}{\( 10x=8x \)}{\( 2x^2+2x=2x^2 \)}{\( x^2+5x=4x+x^2 \)}{}{}}
\End

\Data\ID{1003123801}
\Updated{2020-08-22 04:08:21}
\Level{A}
\SubArea{Quadratic equations and inequalities}
\LANGen{
\Question{
Choose the equation which has no real roots.}
{\( x^2+\sqrt2=0 \)}{\( x^2+\sqrt2 x=0 \)}{\( x^2-\sqrt2=0 \)}{\( x^2-\sqrt2  x=0 \)}{}{}}
\LANGpl{
\Question{
Wybierz równanie, które nie posiada pierwiastków rzeczywistych.}
{\( x^2+\sqrt2=0 \)}{\( x^2+\sqrt2 x=0 \)}{\( x^2-\sqrt2=0 \)}{\( x^2-\sqrt2  x=0 \)}{}{}}
\LANGcs{
\Question{
Vyberte rovnici, která nemá reálné kořeny.}
{\( x^2+\sqrt2=0 \)}{\( x^2+\sqrt2 x=0 \)}{\( x^2-\sqrt2=0 \)}{\( x^2-\sqrt2  x=0 \)}{}{}}
\LANGsk{
\Question{
Vyberte rovnicu, ktorá nemá reálne korene.}
{\( x^2+\sqrt2=0 \)}{\( x^2+\sqrt2 x=0 \)}{\( x^2-\sqrt2=0 \)}{\( x^2-\sqrt2  x=0 \)}{}{}}
\LANGes{
\Question{
Elige la ecuación que no tenga raices reales.}
{\( x^2+\sqrt2=0 \)}{\( x^2+\sqrt2 x=0 \)}{\( x^2-\sqrt2=0 \)}{\( x^2-\sqrt2  x=0 \)}{}{}}
\End

\Data\ID{1003123802}
\Updated{2020-08-22 04:08:39}
\Level{A}
\SubArea{Quadratic equations and inequalities}
\LANGen{
\Question{
Choose the equation which has at least one root in \( (0;1] \).}
{\( 4x^2-3x=0 \)}{\( 3x^2-4x=0 \)}{\( 4x^2+3x=0 \)}{\( 3x^2+4x=0 \)}{}{}}
\LANGpl{
\Question{
Wybierz równanie, które posiada przynajmniej jeden pierwiastek w przedziale \( (0;1\rangle \).}
{\( 4x^2-3x=0 \)}{\( 3x^2-4x=0 \)}{\( 4x^2+3x=0 \)}{\( 3x^2+4x=0 \)}{}{}}
\LANGcs{
\Question{
Vyberte rovnici, jejíž alespoň jeden kořen náleží intervalu \( (0;1\rangle \).}
{\( 4x^2-3x=0 \)}{\( 3x^2-4x=0 \)}{\( 4x^2+3x=0 \)}{\( 3x^2+4x=0 \)}{}{}}
\LANGsk{
\Question{
Vyberte rovnicu, ktorej aspoň jeden koreň patrí do intervalu \( (0;1\rangle \).}
{\( 4x^2-3x=0 \)}{\( 3x^2-4x=0 \)}{\( 4x^2+3x=0 \)}{\( 3x^2+4x=0 \)}{}{}}
\LANGes{
\Question{
Elige la ecuación que tenga por lo menos una raíz en el intervalo \( (0;1] \).}
{\( 4x^2-3x=0 \)}{\( 3x^2-4x=0 \)}{\( 4x^2+3x=0 \)}{\( 3x^2+4x=0 \)}{}{}}
\End

\Data\ID{1003123803}
\Updated{2020-08-22 04:08:00}
\Level{A}
\SubArea{Quadratic equations and inequalities}
\LANGen{
\Question{
Choose the equation which has at least one root in \( (2; 3) \).}
{\( 2x^2-9=0 \)}{\( 2x^2+9=0 \)}{\( 3x^2-9=0 \)}{\( x^2-9=0 \)}{}{}}
\LANGpl{
\Question{
Wybierz równanie, które posiada przynajmniej jeden pierwiastek w przedziale \( (2; 3) \).}
{\( 2x^2-9=0 \)}{\( 2x^2+9=0 \)}{\( 3x^2-9=0 \)}{\( x^2-9=0 \)}{}{}}
\LANGcs{
\Question{
Vyberte rovnici, jejíž alespoň jeden kořen náleží intervalu \( (2; 3) \).}
{\( 2x^2-9=0 \)}{\( 2x^2+9=0 \)}{\( 3x^2-9=0 \)}{\( x^2-9=0 \)}{}{}}
\LANGsk{
\Question{
Vyberte rovnicu, ktorej aspoň jeden koreň patrí do intervalu \( (2; 3) \).}
{\( 2x^2-9=0 \)}{\( 2x^2+9=0 \)}{\( 3x^2-9=0 \)}{\( x^2-9=0 \)}{}{}}
\LANGes{
\Question{
Elige la ecuación que tenga por lo menos una raíz en el intervalo \( (2; 3) \).}
{\( 2x^2-9=0 \)}{\( 2x^2+9=0 \)}{\( 3x^2-9=0 \)}{\( x^2-9=0 \)}{}{}}
\End

\Data\ID{1003123804}
\Updated{2020-08-22 04:08:15}
\Level{A}
\SubArea{Quadratic equations and inequalities}
\LANGen{
\Question{
Choose the equation which has no real roots.}
{\( 4x^2+16=0 \)}{\( -4x^2+16=0 \)}{\( 16x^2-4=0 \)}{\( \sqrt2x^2-2=0 \)}{}{}}
\LANGpl{
\Question{
Wybierz równanie, które nie posiada pierwiastków rzeczywistych.}
{\( 4x^2+16=0 \)}{\( -4x^2+16=0 \)}{\( 16x^2-4=0 \)}{\( \sqrt2x^2-2=0 \)}{}{}}
\LANGcs{
\Question{
Vyberte rovnici, která nemá reálné kořeny.}
{\( 4x^2+16=0 \)}{\( -4x^2+16=0 \)}{\( 16x^2-4=0 \)}{\( \sqrt2x^2-2=0 \)}{}{}}
\LANGsk{
\Question{
Vyberte rovnicu, ktorá nemá reálne korene.}
{\( 4x^2+16=0 \)}{\( -4x^2+16=0 \)}{\( 16x^2-4=0 \)}{\( \sqrt2x^2-2=0 \)}{}{}}
\LANGes{
\Question{
Elige la ecuación que no tenga raíces reales.}
{\( 4x^2+16=0 \)}{\( -4x^2+16=0 \)}{\( 16x^2-4=0 \)}{\( \sqrt2x^2-2=0 \)}{}{}}
\End

\Data\ID{1003123805}
\Updated{2020-08-22 04:08:37}
\Level{A}
\SubArea{Quadratic equations and inequalities}
\LANGen{
\Question{
Choose the equation which has the double root equal to zero.}
{\( 2x^2=0 \)}{\( x^2+x=0 \)}{\( x^2-x=0 \)}{\( -x^2+x=0 \)}{}{}}
\LANGpl{
\Question{
Wybierz równanie, którego podwójny pierwiastek równy jest zeru.}
{\( 2x^2=0 \)}{\( x^2+x=0 \)}{\( x^2-x=0 \)}{\( -x^2+x=0 \)}{}{}}
\LANGcs{
\Question{
Vyberte rovnici, která má jediný, a to nulový kořen.}
{\( 2x^2=0 \)}{\( x^2+x=0 \)}{\( x^2-x=0 \)}{\(-x^2+x=0 \)}{}{}}
\LANGsk{
\Question{
Vyberte rovnicu, ktorá má jediný, a to nulový koreň.}
{\( 2x^2=0 \)}{\( x^2+x=0 \)}{\( x^2-x=0 \)}{\( -x^2+x=0 \)}{}{}}
\LANGes{
\Question{
Elige la ecuación que tiene una única raíz igual a cero.}
{\( 2x^2=0 \)}{\( x^2+x=0 \)}{\( x^2-x=0 \)}{\( -x^2+x=0 \)}{}{}}
\End

\Data\ID{1003123806}
\Updated{2020-08-09 06:08:03}
\Level{A}
\SubArea{Quadratic equations and inequalities}
\LANGen{
\Question{
Choose the interval in which all the roots of the following equation lie.
\[ 5x^2-7x=0 \]}
{\( \left[-\frac75;\frac75\right] \)}{\( \left[-1;1\right] \)}{\( ( 0;2] \)}{\( \left[-\frac75;0\right] \)}{}{}}
\LANGpl{
\Question{
Wybierz przedział, w którym znajdują się wszystkie pierwiastki podanego równania.
\[ 5x^2-7x=0 \]}
{\( \left\langle-\frac75;\frac75\right\rangle \)}{\( \left\langle-1;1\right\rangle \)}{\( ( 0;2\rangle \)}{\( \left\langle-\frac75;0\right\rangle \)}{}{}}
\LANGcs{
\Question{
Vyberte interval, ve kterém leží všechny kořeny dané rovnice.
\[ 5x^2-7x=0 \]}
{\( \left\langle-\frac75;\frac75\right\rangle \)}{\( \left\langle-1;1\right\rangle \)}{\( ( 0;2\rangle \)}{\( \left\langle-\frac75;0\right\rangle \)}{}{}}
\LANGsk{
\Question{
Vyberte interval, v ktorom ležia všetky korene danej rovnice.
\[ 5x^2-7x=0 \]}
{\( \left\langle-\frac75;\frac75\right\rangle \)}{\( \left\langle-1;1\right\rangle \)}{\( ( 0;2\rangle \)}{\( \left\langle-\frac75;0\right\rangle \)}{}{}}
\LANGes{
\Question{
Elige el intervalo en el que están todas las raíces de la siguiente ecuación.
\[ 5x^2-7x=0 \]}
{\( \left[-\frac75;\frac75\right] \)}{\( \left[-1;1\right] \)}{\( ( 0;2] \)}{\( \left[-\frac75;0\right] \)}{}{}}
\End

\Data\ID{1003123807}
\Updated{2020-08-09 06:08:28}
\Level{A}
\SubArea{Quadratic equations and inequalities}
\LANGen{
\Question{
Choose the interval in which at least one root of the following equation lies.
\[ x^2-8=0 \]}
{\( [ 2;3 ] \)}{\( [ 3;4 ] \)}{\( [ -1;1 ] \)}{\( [ -8;-5 ] \)}{}{}}
\LANGpl{
\Question{
Wybierz przedział, w którym znajduje się co najmniej jeden pierwiastek podanego równania. 
\[ x^2-8=0 \]}
{\( \langle 2;3 \rangle \)}{\( \langle 3;4 \rangle \)}{\( \langle -1;1 \rangle \)}{\( \langle -8;-5 \rangle \)}{}{}}
\LANGcs{
\Question{
Vyberte interval, ve kterém leží alespoň jeden kořen dané rovnice.
\[ x^2-8=0 \]}
{\( \langle 2;3 \rangle \)}{\( \langle 3;4 \rangle \)}{\( \langle -1;1 \rangle \)}{\( \langle -8;-5 \rangle \)}{}{}}
\LANGsk{
\Question{
Vyberte interval, v ktorom leží aspoň jeden koreň danej rovnice.
\[ x^2-8=0 \]}
{\( \langle 2;3 \rangle \)}{\( \langle 3;4 \rangle \)}{\( \langle -1;1 \rangle \)}{\( \langle -8;-5 \rangle \)}{}{}}
\LANGes{
\Question{
Elige el intervalo en el que está por lo menos una raíz de la siguiente ecuación. 
\[ x^2-8=0 \]}
{\( [ 2;3 ] \)}{\( [ 3;4 ] \)}{\( [ -1;1 ] \)}{\( [ -8;-5 ] \)}{}{}}
\End

\Data\ID{9000020401}
\Updated{2020-08-09 07:08:24}
\Level{A}
\SubArea{Quadratic equations and inequalities}
\LANGen{
\Question{
Solve the following quadratic equation. 
\[ 
 -x^{2} + 12x - 20 = 0
\]}
{\(x_{1} = 2\),
\(x_{2} = 10\)}{\(x_{1} = -2\),
\(x_{2} = 10\)}{\(x_{1} = -2\),
\(x_{2} = -10\)}{\(x_{1} = 2\),
\(x_{2} = -10\)}{}{}}
\LANGpl{
\Question{
Rozwiąż podane równanie kwadratowe.
\[ 
 -x^{2} + 12x - 20 = 0
\]}
{\(x_{1} = 2\),
\(x_{2} = 10\)}{\(x_{1} = -2\),
\(x_{2} = 10\)}{\(x_{1} = -2\),
\(x_{2} = -10\)}{\(x_{1} = 2\),
\(x_{2} = -10\)}{}{}}
\LANGcs{
\Question{
Přiřaďte kvadratické rovnici \(- x^{2} + 12x - 20 = 0\) 
její kořeny.}
{\(x_1=2\),
\(x_2=10\)}{\(x_1=- 2\),
\(x_2=10\)}{\(x_1=- 2\),
\(x_2=- 10\)}{\(x_1=2\),
\(x_2=- 10\)}{}{}}
\LANGsk{
\Question{
Vyriešte danú kvadratickú rovnicu. 
\[ 
 -x^{2} + 12x - 20 = 0
\]}
{\(x_{1} = 2\),
\(x_{2} = 10\)}{\(x_{1} = -2\),
\(x_{2} = 10\)}{\(x_{1} = -2\),
\(x_{2} = -10\)}{\(x_{1} = 2\),
\(x_{2} = -10\)}{}{}}
\LANGes{
\Question{
Resuelve la siguiente ecuación cuadrática.  
\[ 
 -x^{2} + 12x - 20 = 0
\]}
{\(x_{1} = 2\),
\(x_{2} = 10\)}{\(x_{1} = -2\),
\(x_{2} = 10\)}{\(x_{1} = -2\),
\(x_{2} = -10\)}{\(x_{1} = 2\),
\(x_{2} = -10\)}{}{}}
\End

\Data\ID{9000020402}
\Updated{2020-08-20 11:08:06}
\Level{A}
\SubArea{Quadratic equations and inequalities}
\LANGen{
\Question{
Identify an equation which does not have real solution.}
{\(x^{2} - 2x + 5 = 0\)}{\(x^{2} - 5 = 0\)}{\(x^{2} + 0.8x = 0\)}{\(- x^{2} + 2x + 35 = 0\)}{}{}}
\LANGpl{
\Question{
Określ, które równanie nie ma rzeczywistego rozwiązania.}
{\(x^{2} - 2x + 5 = 0\)}{\(x^{2} - 5 = 0\)}{\(x^{2} + 0.8x = 0\)}{\(- x^{2} + 2x + 35 = 0\)}{}{}}
\LANGcs{
\Question{
Určete, která z uvedených rovnic nemá reálné kořeny.}
{\(x^{2} - 2x + 5 = 0\)}{\(x^{2} - 5 = 0\)}{\(x^{2} + 0{,}8x = 0\)}{\(- x^{2} + 2x + 35 = 0\)}{}{}}
\LANGsk{
\Question{
Určte, ktorá z uvedených rovníc nemá reálne korene.}
{\(x^{2} - 2x + 5 = 0\)}{\(x^{2} - 5 = 0\)}{\(x^{2} + 0.8x = 0\)}{\(- x^{2} + 2x + 35 = 0\)}{}{}}
\LANGes{
\Question{
Identifica la ecuación que no tiene solución real.}
{\(x^{2} - 2x + 5 = 0\)}{\(x^{2} - 5 = 0\)}{\(x^{2} + 0.8x = 0\)}{\(- x^{2} + 2x + 35 = 0\)}{}{}}
\End

\Data\ID{9000020403}
\Updated{2020-08-20 11:08:56}
\Level{A}
\SubArea{Quadratic equations and inequalities}
\LANGen{
\Question{
Identify an equation which does not have at least one solution in the interval
\((0;\infty )\).}
{\(x^{2} + 5x + 6 = 0\)}{\(x^{2} - 2x - 3 = 0\)}{\(x^{2} - 10x = 0\)}{\(x^{2} - 10x + 24 = 0\)}{}{}}
\LANGpl{
\Question{
Określ, które równanie nie ma przynajmniej jednego rozwiązania w przedziale
\((0;\infty )\).}
{\(x^{2} + 5x + 6 = 0\)}{\(x^{2} - 2x - 3 = 0\)}{\(x^{2} - 10x = 0\)}{\(x^{2} - 10x + 24 = 0\)}{}{}}
\LANGcs{
\Question{
Určete, která z uvedených rovnic nemá aspoň jeden kořen z intervalu
\((0;\infty )\).}
{\(x^{2} + 5x + 6 = 0\)}{\(x^{2} - 2x - 3 = 0\)}{\(x^{2} - 10x = 0\)}{\(x^{2} - 10x + 24 = 0\)}{}{}}
\LANGsk{
\Question{
Určte, ktorá z uvedených rovníc nemá aspoň jeden koreň z intervalu
\((0;\infty )\).}
{\(x^{2} + 5x + 6 = 0\)}{\(x^{2} - 2x - 3 = 0\)}{\(x^{2} - 10x = 0\)}{\(x^{2} - 10x + 24 = 0\)}{}{}}
\LANGes{
\Question{
Identifica la ecuación que no tiene por lo menos una solución en el intervalo 
\((0;\infty )\).}
{\(x^{2} + 5x + 6 = 0\)}{\(x^{2} - 2x - 3 = 0\)}{\(x^{2} - 10x = 0\)}{\(x^{2} - 10x + 24 = 0\)}{}{}}
\End

\Data\ID{9000020404}
\Updated{2020-08-22 04:08:03}
\Level{A}
\SubArea{Quadratic equations and inequalities}
\LANGen{
\Question{
Find a number which is a sum of one half of the bigger solution of 
\[ 
 x^{2} - 10x + 24 = 0
\]
and a double of the smaller solution of 
\[ 
 -x^{2} + 10x - 16 = 0.
\]}
{\(7\)}{\(12\)}{\(6\)}{\(14\)}{}{}}
\LANGpl{
\Question{
Znajdź liczbę, która jest sumą jednej drugiej większego rozwiązania równania
\[ 
 x^{2} - 10x + 24 = 0
\]
i dwukrotności mniejszego rozwiązania równania
\[ 
 -x^{2} + 10x - 16 = 0.
\]}
{\(7\)}{\(12\)}{\(6\)}{\(14\)}{}{}}
\LANGcs{
\Question{
Jaké získáme číslo, jestliže sečteme polovinu většího kořenu rovnice
\[x^{2} - 10x + 24 = 0\] a dvojnásobek menšího
kořenu rovnice \[- x^{2} + 10x - 16 = 0?\]}
{\(7\)}{\(12\)}{\(6\)}{\(14\)}{}{}}
\LANGsk{
\Question{
Nájdi číslo, ktoré získame, ak sčítame polovicu väčšieho koreňa rovnice  
\[ 
 x^{2} - 10x + 24 = 0
\]
a dvojnásobok menšieho koreňa rovnice 
\[ 
 -x^{2} + 10x - 16 = 0.
\]}
{\(7\)}{\(12\)}{\(6\)}{\(14\)}{}{}}
\LANGes{
\Question{
Halla el número que es suma de la mitad de la solución más alta de la ecuación  
\[ 
 x^{2} - 10x + 24 = 0
\]
y el doble de la solución más baja de la ecuación  
\[ 
 -x^{2} + 10x - 16 = 0.
\]}
{\(7\)}{\(12\)}{\(6\)}{\(14\)}{}{}}
\End

\Data\ID{9000020405}
\Updated{2020-08-20 07:08:11}
\Level{A}
\SubArea{Quadratic equations and inequalities}
\LANGen{
\Question{
Identify an equation which does not have the set
\(K = \{ - 3;6\}\) as the
set of all solutions of this equation.}
{\(3x^{2} - 9x + 54 = 0\)}{\(2x^{2} - 6x - 36 = 0\)}{\(\frac{1}
{3}x^{2} - x - 6 = 0\)}{\(- x^{2} + 3x + 18 = 0\)}{}{}}
\LANGpl{
\Question{
Określ, dla którego równania zbiór
\(K = \{ - 3;6\}\) nie jest zbiorem rozwiązań.}
{\(3x^{2} - 9x + 54 = 0\)}{\(2x^{2} - 6x - 36 = 0\)}{\(\frac{1}
{3}x^{2} - x - 6 = 0\)}{\(- x^{2} + 3x + 18 = 0\)}{}{}}
\LANGcs{
\Question{
Určete, která z uvedených rovnic nemá množinu kořenů
\(K = \{ - 3;6\}\).}
{\(3x^{2} - 9x + 54 = 0\)}{\(2x^{2} - 6x - 36 = 0\)}{\(\frac{1}
{3}x^{2} - x - 6 = 0\)}{\(- x^{2} + 3x + 18 = 0\)}{}{}}
\LANGsk{
\Question{
Určte, ktorá z uvedených rovníc nemá množinu koreňov
\(K = \{ - 3;6\}\).}
{\(3x^{2} - 9x + 54 = 0\)}{\(2x^{2} - 6x - 36 = 0\)}{\(\frac{1}
{3}x^{2} - x - 6 = 0\)}{\(- x^{2} + 3x + 18 = 0\)}{}{}}
\LANGes{
\Question{
Identifica la ecuación que no tiene el conjunto de soluciones \(K = \{ - 3;6\}\).}
{\(3x^{2} - 9x + 54 = 0\)}{\(2x^{2} - 6x - 36 = 0\)}{\(\frac{1}
{3}x^{2} - x - 6 = 0\)}{\(- x^{2} + 3x + 18 = 0\)}{}{}}
\End

\Data\ID{9000020407}
\Updated{2020-08-20 08:08:09}
\Level{A}
\SubArea{Quadratic equations and inequalities}
\LANGen{
\Question{
In the following list identify an equation with real solution.}
{\(- 0.5x^{2} + 2x + 3 = 0\)}{\(- x^{2} + 4x - 5 = 0\)}{\(2x^{2} - 3x + 3 = 0\)}{\(x^{2} - x + 1 = 0\)}{}{}}
\LANGpl{
\Question{
Które równanie z niżej podanych ma rzeczywiste rozwiązanie?}
{\(- 0.5x^{2} + 2x + 3 = 0\)}{\(- x^{2} + 4x - 5 = 0\)}{\(2x^{2} - 3x + 3 = 0\)}{\(x^{2} - x + 1 = 0\)}{}{}}
\LANGcs{
\Question{
Vyberte z uvedených rovnic takovou, která má reálné kořeny.}
{\(- 0{,}5x^{2} + 2x + 3 = 0\)}{\(- x^{2} + 4x - 5 = 0\)}{\(2x^{2} - 3x + 3 = 0\)}{\(x^{2} - x + 1 = 0\)}{}{}}
\LANGsk{
\Question{
Vyberte z uvedených rovníc takú, ktorá má reálne korene.}
{\(- 0.5x^{2} + 2x + 3 = 0\)}{\(- x^{2} + 4x - 5 = 0\)}{\(2x^{2} - 3x + 3 = 0\)}{\(x^{2} - x + 1 = 0\)}{}{}}
\LANGes{
\Question{
Identifica de la siguiente lista la ecuación con soluciones reales.}
{\(- 0.5x^{2} + 2x + 3 = 0\)}{\(- x^{2} + 4x - 5 = 0\)}{\(2x^{2} - 3x + 3 = 0\)}{\(x^{2} - x + 1 = 0\)}{}{}}
\End

\Data\ID{9000020408}
\Updated{2020-08-20 08:08:23}
\Level{A}
\SubArea{Quadratic equations and inequalities}
\LANGen{
\Question{
Which of the given equations have at least one root the same?
\[ \begin{aligned}
 x^{2} + 8x + 15 & = 0 &\text{(1)}
 \\x^{2} - 8x + 15 & = 0 &\text{(2)}
 \\x^{2} +\phantom{ 8}x - 12 & = 0 &\text{(3)}
 \\x^{2} - 2x -\phantom{ 1}8 & = 0 &\text{(4)}
\\\end{aligned}\]}
{equations (2) and (3)}{equations (1) and (3)}{equations (2) and (4)}{Such a pair does not exist.}{}{}}
\LANGpl{
\Question{
Spośród podanych równań kwadratowych wybierz parę równań, która ma wspólny pierwiastek.
\[ \begin{aligned}
 x^{2} + 8x + 15 & = 0 &\text{(1)}
 \\x^{2} - 8x + 15 & = 0 &\text{(2)}
 \\x^{2} +\phantom{ 8}x - 12 & = 0 &\text{(3)}
 \\x^{2} - 2x -\phantom{ 1}8 & = 0 &\text{(4)}
\\\end{aligned}\]}
{równania (2) i (3)}{równania (1) i (3)}{równania (2) i (4)}{Nie ma takiej pary równań.}{}{}}
\LANGcs{
\Question{
Uveďte, které z uvedených rovnic mají společný kořen.
\[ \begin{aligned}
 x^{2} + 8x + 15 & = 0 &\text{(1)}
 \\x^{2} - 8x + 15 & = 0 &\text{(2)}
 \\x^{2} +\phantom{ 8}x - 12 & = 0 &\text{(3)}
 \\x^{2} - 2x -\phantom{ 1}8 & = 0 &\text{(4)}
\\\end{aligned}\]}
{rovnice (2) a (3)}{rovnice (1) a (3)}{rovnice (2) a (4)}{Žádné dvě z uvedených rovnic nemají společný kořen.}{}{}}
\LANGsk{
\Question{
Zistite, ktoré z uvedených rovníc majú spoločný koreň.
\[ \begin{aligned}
 x^{2} + 8x + 15 & = 0 &\text{(1)}
 \\x^{2} - 8x + 15 & = 0 &\text{(2)}
 \\x^{2} +\phantom{ 8}x - 12 & = 0 &\text{(3)}
 \\x^{2} - 2x -\phantom{ 1}8 & = 0 &\text{(4)}
\\\end{aligned}\]}
{rovnice (2) a (3)}{rovnice (1) a (3)}{rovnice (2) a (4)}{Žiadne dve z uvedených nemajú spoločný koreň.}{}{}}
\LANGes{
\Question{
¿Cuáles de las ecuaciones dadas tienen por lo menos una raíz en común? 
\[ \begin{aligned}
 x^{2} + 8x + 15 & = 0 &\text{(1)}
 \\x^{2} - 8x + 15 & = 0 &\text{(2)}
 \\x^{2} +\phantom{ 8}x - 12 & = 0 &\text{(3)}
 \\x^{2} - 2x -\phantom{ 1}8 & = 0 &\text{(4)}
\\\end{aligned}\]}
{ecuaciones (2) y (3)}{ecuaciones (1) y (3)}{ecuaciones (2) y (4)}{Esta pareja no existe.}{}{}}
\End

\Data\ID{9000025601}
\Updated{2020-08-21 05:08:03}
\Level{A}
\SubArea{Quadratic equations and inequalities}
\LANGen{
\Question{
Which of the quadratic equations from the following list has all solutions in the interval
\([ - 5;3] \)?}
{\(x^{2} - 2x - 3 = 0\)}{\(x^{2} - 3x - 10 = 0\)}{\(x^{2} - 2x - 15 = 0\)}{\(x^{2} - 3x - 18 = 0\)}{}{}}
\LANGpl{
\Question{
Które z podanych poniżej równań kwadratowych ma swoje rozwiązania w przedziale \([ - 5;3] \)?}
{\(x^{2} - 2x - 3 = 0\)}{\(x^{2} - 3x - 10 = 0\)}{\(x^{2} - 2x - 15 = 0\)}{\(x^{2} - 3x - 18 = 0\)}{}{}}
\LANGcs{
\Question{
Která z následujících kvadratických rovnic má všechny kořeny v intervalu
\(\langle - 5;3\rangle \)?}
{\(x^{2} - 2x - 3 = 0\)}{\(x^{2} - 3x - 10 = 0\)}{\(x^{2} - 2x - 15 = 0\)}{\(x^{2} - 3x - 18 = 0\)}{}{}}
\LANGsk{
\Question{
Ktorá z nasledujúcich kvadratických rovníc má všetky korene v intervale
\([ - 5;3] \)?}
{\(x^{2} - 2x - 3 = 0\)}{\(x^{2} - 3x - 10 = 0\)}{\(x^{2} - 2x - 15 = 0\)}{\(x^{2} - 3x - 18 = 0\)}{}{}}
\LANGes{
\Question{
¿Cuál de las siguientes ecuaciones tiene todas las soluciones en el intervalo 
\([ - 5;3] \)?}
{\(x^{2} - 2x - 3 = 0\)}{\(x^{2} - 3x - 10 = 0\)}{\(x^{2} - 2x - 15 = 0\)}{\(x^{2} - 3x - 18 = 0\)}{}{}}
\End

\Data\ID{9000025602}
\Updated{2020-08-21 06:08:22}
\Level{A}
\SubArea{Quadratic equations and inequalities}
\LANGen{
\Question{
In the following list identify a set which contains at least one solution of the quadratic
equation \(x^{2} - 121 = 0\).}
{\(\{ - 11;1;13\}\)}{\(\{ - 5;0;5;10\}\)}{\(\{3;7;9;19\}\)}{\(\{ - 15;-12;-7\}\)}{}{}}
\LANGpl{
\Question{
Z podanej listy wybierz zbiór, który zawiera przynajmniej jedno rozwiązanie podanego równania kwadratowego  \(x^{2} - 121 = 0\).}
{\(\{ - 11;1;13\}\)}{\(\{ - 5;0;5;10\}\)}{\(\{3;7;9;19\}\)}{\(\{ - 15;-12;-7\}\)}{}{}}
\LANGcs{
\Question{
Vyberte množinu, ve které se nachází aspoň jeden z kořenů kvadratické
rovnice \(x^{2} - 121 = 0\).}
{\(\{ - 11;1;13\}\)}{\(\{ - 5;0;5;10\}\)}{\(\{3;7;9;19\}\)}{\(\{ - 15;-12;-7\}\)}{}{}}
\LANGsk{
\Question{
Vyberte množinu, v ktorej sa nachádza aspoň jeden z koreňov kvadratickej rovnice \(x^{2} - 121 = 0\).}
{\(\{ - 11;1;13\}\)}{\(\{ - 5;0;5;10\}\)}{\(\{3;7;9;19\}\)}{\(\{ - 15;-12;-7\}\)}{}{}}
\LANGes{
\Question{
Identifica el conjunto en el que se encuentra por lo menos una solución de la siguiente ecuación cuadrática \(x^{2} - 121 = 0\).}
{\(\{ - 11;1;13\}\)}{\(\{ - 5;0;5;10\}\)}{\(\{3;7;9;19\}\)}{\(\{ - 15;-12;-7\}\)}{}{}}
\End

\Data\ID{9000025603}
\Updated{2020-08-20 09:08:11}
\Level{A}
\SubArea{Quadratic equations and inequalities}
\LANGen{
\Question{
In the following list identify an equation equivalent to the following quadratic
equation. 
\[ 
 2(x - 8)\left (x + \frac{1} 
{2}\right ) = 0
\]}
{\(2x^{2} - 15x - 8 = 0\)}{\(2x^{2} - 8x + \frac{1} 
{2} = 0\)}{\(8x^{2} - 15x + \frac{1} 
{2} = 0\)}{\(8x^{2} - 8x - 1 = 0\)}{}{}}
\LANGpl{
\Question{
Z podanej listy wybierz równanie, które jest równoważne podanemu równaniu kwadratowemu: 
\[ 
 2(x - 8)\left (x + \frac{1} 
{2}\right ) = 0
\]}
{\(2x^{2} - 15x - 8 = 0\)}{\(2x^{2} - 8x + \frac{1} 
{2} = 0\)}{\(8x^{2} - 15x + \frac{1} 
{2} = 0\)}{\(8x^{2} - 8x - 1 = 0\)}{}{}}
\LANGcs{
\Question{
Která z kvadratických rovnic je ekvivalentní rovnici
\(2(x - 8)\left (x + \frac{1} 
{2}\right ) = 0\)?}
{\(2x^{2} - 15x - 8 = 0\)}{\(2x^{2} - 8x + \frac{1} 
{2} = 0\)}{\(8x^{2} - 15x + \frac{1} 
{2} = 0\)}{\(8x^{2} - 8x - 1 = 0\)}{}{}}
\LANGsk{
\Question{
Ktorá z kvadratických rovníc je ekvivalentná danej rovnici? 
\[ 
 2(x - 8)\left (x + \frac{1} 
{2}\right ) = 0
\]}
{\(2x^{2} - 15x - 8 = 0\)}{\(2x^{2} - 8x + \frac{1} 
{2} = 0\)}{\(8x^{2} - 15x + \frac{1} 
{2} = 0\)}{\(8x^{2} - 8x - 1 = 0\)}{}{}}
\LANGes{
\Question{
¿Cuál de las siguientes ecuaciones cuadráticas equivale a la siguiente ecuación?  
\[ 
 2(x - 8)\left (x + \frac{1} 
{2}\right ) = 0
\]}
{\(2x^{2} - 15x - 8 = 0\)}{\(2x^{2} - 8x + \frac{1} 
{2} = 0\)}{\(8x^{2} - 15x + \frac{1} 
{2} = 0\)}{\(8x^{2} - 8x - 1 = 0\)}{}{}}
\End

\Data\ID{9000025604}
\Updated{2020-08-21 07:08:30}
\Level{A}
\SubArea{Quadratic equations and inequalities}
\LANGen{
\Question{
In the following list identify an equation which does not have solution in the set of
real numbers.}
{\(8x^{2} - x + 1 = 0\)}{\(8x^{2} + 8x - 1 = 0\)}{\(8x^{2} - 8x + 1 = 0\)}{\(8x^{2} - x - 1 = 0\)}{}{}}
\LANGpl{
\Question{
Które z podanych poniżej równań nie ma rozwiązania w zbiorze liczb rzeczywistych?}
{\(8x^{2} - x + 1 = 0\)}{\(8x^{2} + 8x - 1 = 0\)}{\(8x^{2} - 8x + 1 = 0\)}{\(8x^{2} - x - 1 = 0\)}{}{}}
\LANGcs{
\Question{
Která z kvadratických rovnic nemá řešení v množině
\(\mathbb{R}\)?}
{\(8x^{2} - x + 1 = 0\)}{\(8x^{2} + 8x - 1 = 0\)}{\(8x^{2} - 8x + 1 = 0\)}{\(8x^{2} - x - 1 = 0\)}{}{}}
\LANGsk{
\Question{
Ktorá z kvadratických rovníc nemá riešenie v množine reálnych čísel?}
{\(8x^{2} - x + 1 = 0\)}{\(8x^{2} + 8x - 1 = 0\)}{\(8x^{2} - 8x + 1 = 0\)}{\(8x^{2} - x - 1 = 0\)}{}{}}
\LANGes{
\Question{
¿Cuál de las siguientes ecuaciones no tiene solución en el conjunto de números reales?}
{\(8x^{2} - x + 1 = 0\)}{\(8x^{2} + 8x - 1 = 0\)}{\(8x^{2} - 8x + 1 = 0\)}{\(8x^{2} - x - 1 = 0\)}{}{}}
\End

\Data\ID{9000025605}
\Updated{2020-08-21 07:08:49}
\Level{A}
\SubArea{Quadratic equations and inequalities}
\LANGen{
\Question{
Find the interval which contains all the solutions of the following quadratic
equation.
\[ 
 \text{$5x^{2} - 3x - 2 = 0$}
\]}
{\((-0.5;2)\)}{\([ - 1;0] \)}{\([ 0;4)\)}{\([ - 3;1)\)}{}{}}
\LANGpl{
\Question{
Znajdź przedział, który zawiera wszystkie rozwiązania następującego równania kwadratowego: 
\[ 
 \text{$5x^{2} - 3x - 2 = 0$}
\]}
{\((-0.5;2)\)}{\([ - 1;0] \)}{\([ 0;4)\)}{\([ - 3;1)\)}{}{}}
\LANGcs{
\Question{
Vyberte interval, ve kterém se nachází všechny kořeny kvadratické rovnice
\(5x^{2} - 3x - 2 = 0\).}
{\(\left (-0{,}5;2\right )\)}{\(\left \langle -1;0\right \rangle \)}{\(\left \langle 0;4\right )\)}{\(\left \langle -3;1\right )\)}{}{}}
\LANGsk{
\Question{
Vyberte interval, v ktorom sa nachádzajú všetky korene kvadratickej rovnice.  
\[ 
 \text{$5x^{2} - 3x - 2 = 0$}
\]}
{\((-0.5;2)\)}{\([ - 1;0] \)}{\([ 0;4)\)}{\([ - 3;1)\)}{}{}}
\LANGes{
\Question{
Halla el intervalo que contiene todas las soluciones de la siguiente ecuación cuadrática. 
\[ 
 \text{$5x^{2} - 3x - 2 = 0$}
\]}
{\((-0.5;2)\)}{\([ - 1;0] \)}{\([ 0;4)\)}{\([ - 3;1)\)}{}{}}
\End

\Data\ID{9000025606}
\Updated{2020-08-22 05:08:13}
\Level{A}
\SubArea{Quadratic equations and inequalities}
\LANGen{
\Question{
In the following list identify a quadratic equation which has a unique solution.}
{\(x^{2} + 2x + 1 = 0\)}{\(x^{2} - 3x - 1 = 0\)}{\(x^{2} + 2x - 1 = 0\)}{\(x^{2} - 3x + 1 = 0\)}{}{}}
\LANGpl{
\Question{
Z podanej listy wybierz równanie kwadratowe, które ma tylko jedno rozwiązanie.}
{\(x^{2} + 2x + 1 = 0\)}{\(x^{2} - 3x - 1 = 0\)}{\(x^{2} + 2x - 1 = 0\)}{\(x^{2} - 3x + 1 = 0\)}{}{}}
\LANGcs{
\Question{
Která z kvadratických rovnic má právě jedno řešení?}
{\(x^{2} + 2x + 1 = 0\)}{\(x^{2} - 3x - 1 = 0\)}{\(x^{2} + 2x - 1 = 0\)}{\(x^{2} - 3x + 1 = 0\)}{}{}}
\LANGsk{
\Question{
Ktorá z kvadratických rovníc má práve jedno riešenie?}
{\(x^{2} + 2x + 1 = 0\)}{\(x^{2} - 3x - 1 = 0\)}{\(x^{2} + 2x - 1 = 0\)}{\(x^{2} - 3x + 1 = 0\)}{}{}}
\LANGes{
\Question{
Identifica la ecuación cuadrática con una única solución.}
{\(x^{2} + 2x + 1 = 0\)}{\(x^{2} - 3x - 1 = 0\)}{\(x^{2} + 2x - 1 = 0\)}{\(x^{2} - 3x + 1 = 0\)}{}{}}
\End

\Data\ID{9000025607}
\Updated{2020-08-22 05:08:46}
\Level{A}
\SubArea{Quadratic equations and inequalities}
\LANGen{
\Question{
In the following list identify a quadratic equation which cannot be factorized in
\(\mathbb{R}\).}
{\(- 7x^{2} + 3x - 1 = 0\)}{\(5x^{2} + 2x - 3 = 0\)}{\(- 3x^{2} + 2x + 3 = 0\)}{\(6x^{2} + 3x - 1 = 0\)}{}{}}
\LANGpl{
\Question{
Którego z podanych równań kwadratowych nie możemy rozłożyć na czynniki pierwsze w  
\(\mathbb{R}\)?}
{\(- 7x^{2} + 3x - 1 = 0\)}{\(5x^{2} + 2x - 3 = 0\)}{\(- 3x^{2} + 2x + 3 = 0\)}{\(6x^{2} + 3x - 1 = 0\)}{}{}}
\LANGcs{
\Question{
Kterou z kvadratických rovnic nelze v
\(\mathbb{R}\) 
rozložit na součin kořenových činitelů?}
{\(- 7x^{2} + 3x - 1 = 0\)}{\(5x^{2} + 2x - 3 = 0\)}{\(- 3x^{2} + 2x + 3 = 0\)}{\(6x^{2} + 3x - 1 = 0\)}{}{}}
\LANGsk{
\Question{
Ktorú z kvadratických rovníc nemožno v 
\(\mathbb{R}\) rozložiť na súčin koreňových činiteľov?}
{\(- 7x^{2} + 3x - 1 = 0\)}{\(5x^{2} + 2x - 3 = 0\)}{\(- 3x^{2} + 2x + 3 = 0\)}{\(6x^{2} + 3x - 1 = 0\)}{}{}}
\LANGes{
\Question{
Identifica la ecuación cuadrática que no puede descomponerse en factores en 
\(\mathbb{R}\).}
{\(- 7x^{2} + 3x - 1 = 0\)}{\(5x^{2} + 2x - 3 = 0\)}{\(- 3x^{2} + 2x + 3 = 0\)}{\(6x^{2} + 3x - 1 = 0\)}{}{}}
\End

\Data\ID{9000025609}
\Updated{2020-08-22 05:08:16}
\Level{A}
\SubArea{Quadratic equations and inequalities}
\LANGen{
\Question{
In the following list identify an equation such that one of the solutions equals
\(- 1\).}
{\(5x^{2} + 2x - 3 = 0\)}{\(5x^{2} - 2x - 3 = 0\)}{\(5x^{2} + 2x + 3 = 0\)}{\(5x^{2} - 2x + 3 = 0\)}{}{}}
\LANGpl{
\Question{
Z podanych równań wybierz to, dla którego jedno z rozwiązań jest równe 
\(- 1\).}
{\(5x^{2} + 2x - 3 = 0\)}{\(5x^{2} - 2x - 3 = 0\)}{\(5x^{2} + 2x + 3 = 0\)}{\(5x^{2} - 2x + 3 = 0\)}{}{}}
\LANGcs{
\Question{
Která z kvadratických rovnic má jeden z kořenů roven
\(- 1\)?}
{\(5x^{2} + 2x - 3 = 0\)}{\(5x^{2} - 2x - 3 = 0\)}{\(5x^{2} + 2x + 3 = 0\)}{\(5x^{2} - 2x + 3 = 0\)}{}{}}
\LANGsk{
\Question{
Ktorá z kvadratických rovníc má jeden z koreňov rovný
\(- 1\)?}
{\(5x^{2} + 2x - 3 = 0\)}{\(5x^{2} - 2x - 3 = 0\)}{\(5x^{2} + 2x + 3 = 0\)}{\(5x^{2} - 2x + 3 = 0\)}{}{}}
\LANGes{
\Question{
¿Cuál de las siguientes ecuaciones tiene una de las soluciones igual a 
\(- 1\).}
{\(5x^{2} + 2x - 3 = 0\)}{\(5x^{2} - 2x - 3 = 0\)}{\(5x^{2} + 2x + 3 = 0\)}{\(5x^{2} - 2x + 3 = 0\)}{}{}}
\End

\Data\ID{1003032504}
\Updated{2020-08-22 04:08:01}
\Level{B}
\SubArea{Quadratic equations and inequalities}
\LANGen{
\Question{
Factoring the polynomial \( -4x^4+26x^3-12x^2 \) you get:}
{\( -2x^2(x-6)(2x-1) \)}{\( 2x^2(x+1)(2x-1) \)}{\( -2x^2(x+6)(2x+1) \)}{\( 2x^2(x-6)(2x+1) \)}{}{}}
\LANGpl{
\Question{
Po rozłożeniu wielomianu  \( -4x^4+26x^3-12x^2 \) na czynniki otrzymamy:}
{\( -2x^2(x-6)(2x-1) \)}{\( 2x^2(x+1)(2x-1) \)}{\( -2x^2(x+6)(2x+1) \)}{\( 2x^2(x-6)(2x+1) \)}{}{}}
\LANGcs{
\Question{
Rozložením mnohočlenu \( -4x^4+26x^3-12x^2 \) dostaneme:}
{\( -2x^2(x-6)(2x-1) \)}{\( 2x^2(x+1)(2x-1) \)}{\( -2x^2(x+6)(2x+1) \)}{\( 2x^2(x-6)(2x+1) \)}{}{}}
\LANGsk{
\Question{
Rozkladom polynómu \( -4x^4+26x^3-12x^2 \) na súčin dostaneme:}
{\( -2x^2(x-6)(2x-1) \)}{\( 2x^2(x+1)(2x-1) \)}{\( -2x^2(x+6)(2x+1) \)}{\( 2x^2(x-6)(2x+1) \)}{}{}}
\LANGes{
\Question{
Factorizando el polinomio \( -4x^4+26x^3-12x^2 \) obtenemos:}
{\( -2x^2(x-6)(2x-1) \)}{\( 2x^2(x+1)(2x-1) \)}{\( -2x^2(x+6)(2x+1) \)}{\( 2x^2(x-6)(2x+1) \)}{}{}}
\End

\Data\ID{1003085401}
\Updated{2020-08-20 07:08:46}
\Level{B}
\SubArea{Quadratic equations and inequalities}
\LANGen{
\Question{
On their birthday students bring candies for their classmates. The birthday person gives a candy to every other student, not to himself or herself. During a year, \( 650 \) candies have been given away. Find out how many students are there in the class. (Note: All students’ birthdays were on school days.)}
{\( 26 \)}{\( 25 \)}{\( 27 \)}{\( 24 \)}{}{}}
\LANGpl{
\Question{
Z okazji urodzin uczniowie przynoszą cukierki dla swoich kolegów z klasy. Solenizant daje cukierek każdemu uczniowi oprócz siebie. W ciągu roku rozdano, \( 650 \) cukierków. Wyznacz liczbę wszystkich uczniów tej klasy. (Wskazówka: Urodziny wszystkich uczniów miały miejsce w dni szkolne.)}
{\( 26 \)}{\( 25 \)}{\( 27 \)}{\( 24 \)}{}{}}
\LANGcs{
\Question{
Ve třídě rozdávají žáci vždy o svých narozeninách spolužákům bonbóny. Oslavenec dá vždy každému po jednom bonbónu, sobě nedává. Za rok se ve třídě rozdalo celkem \( 650 \) bonbónů. Kolik žáků je ve třídě? (Poznámka: Všichni žáci třídy měli narozeniny v den, kdy se konala výuka.)}
{\( 26 \)}{\( 25 \)}{\( 27 \)}{\( 24 \)}{}{}}
\LANGsk{
\Question{
V deň svojich narodenín študenti rozdávajú v triede svojim spolužiakom bonbóny. Oslávenec dá vždy každému spolužiakovi jeden bonbón, sebe nedáva. Za rok sa v triede rozdalo celkom \( 650 \) bonbónov. Koľko študentov je v triede? (Poznámka: Všetci študenti v triede mali narodeniny v deň, keď sa konalo vyučovanie.)}
{\( 26 \)}{\( 25 \)}{\( 27 \)}{\( 24 \)}{}{}}
\LANGes{
\Question{
En la clase el día de cumpleaños se suelen repartir caramelos. El que celebra su cumple regala a cada uno de sus compañeros, excepto a sí mismo, un caramelo. A lo largo del año se repartieron \( 650 \) caramelos. ¿Cuántos estudiantes hay en la clase? (Nota: Todos los estudiantes celebran su cumpleaños durante el año escolar.)}
{\( 26 \)}{\( 25 \)}{\( 27 \)}{\( 24 \)}{}{}}
\End

\Data\ID{1003085402}
\Updated{2020-08-20 07:08:16}
\Level{B}
\SubArea{Quadratic equations and inequalities}
\LANGen{
\Question{
The product of two consecutive  even numbers is \( 2 808 \). Find their sum.}
{\( 106 \)}{\( 104 \)}{\( 102 \)}{\( 108 \)}{}{}}
\LANGpl{
\Question{
Iloczyn dwóch kolejnych liczb parzystych jest równy \( 2 808 \). Wyznacz ich sumę.}
{\( 106 \)}{\( 104 \)}{\( 102 \)}{\( 108 \)}{}{}}
\LANGcs{
\Question{
Součin dvou po sobě jdoucích sudých čísel je  \( 2 808 \). Určete součet těchto čísel.}
{\( 106 \)}{\( 104 \)}{\( 102 \)}{\( 108 \)}{}{}}
\LANGsk{
\Question{
Súčin dvoch po sebe idúcich párnych čísel je \( 2 808 \). Určte súčet týchto čísel.}
{\( 106 \)}{\( 104 \)}{\( 102 \)}{\( 108 \)}{}{}}
\LANGes{
\Question{
El producto de dos números pares consecutivos es \( 2 808 \). Encuentra su suma.}
{\( 106 \)}{\( 104 \)}{\( 102 \)}{\( 108 \)}{}{}}
\End

\Data\ID{1003085407}
\Updated{2020-08-22 04:08:34}
\Level{B}
\SubArea{Quadratic equations and inequalities}
\LANGen{
\Question{
If the area of the square is increased by \( 64\% \), by what percentage the length of its side is increased? Round the result to an integer.}
{\( 28 \)}{\( 8 \)}{\( 32 \)}{\( 16 \)}{}{}}
\LANGpl{
\Question{
Jeżeli pole kwadratu zwiększymy o  \( 64\% \), o ile procent wzrośnie długość jego boku? Zaokrąglij do liczby całkowitej.}
{\( 28 \)}{\( 8 \)}{\( 32 \)}{\( 16 \)}{}{}}
\LANGcs{
\Question{
Pokud zvětšíme obsah čtverce o \( 64\,\% \), o kolik procent se zvětší délka jeho strany? Výsledek zaokrouhlete na celé číslo.}
{\( 28 \)}{\( 8 \)}{\( 32 \)}{\( 16 \)}{}{}}
\LANGsk{
\Question{
Ak zväčšíme obsah štvorca o \( 64\% \), o koľko percent sa zväčší dĺžka jeho strany? Výsledok zaokrúhlite na celé číslo.}
{\( 28 \)}{\( 8 \)}{\( 32 \)}{\( 16 \)}{}{}}
\LANGes{
\Question{
Si aumentamos el área de un cuadro en un \( 64\% \), ¿En qué porcentaje aumenta su lado? Aproxima el resultado a un número entero.}
{\( 28 \)}{\( 8 \)}{\( 32 \)}{\( 16 \)}{}{}}
\End

\Data\ID{1003085409}
\Updated{2020-08-22 04:08:23}
\Level{B}
\SubArea{Quadratic equations and inequalities}
\LANGen{
\Question{
The area of the rectangle is \( 1\,269\,\mathrm{cm}^2 \). Its length is by \( 20\,\mathrm{cm} \) longer than its width. Find the sum of the length and the width of the rectangle.}
{\( 74\,\mathrm{cm} \)}{\( 35\,\mathrm{cm} \)}{\( 27\,\mathrm{cm} \)}{\( 57\,\mathrm{cm} \)}{}{}}
\LANGpl{
\Question{
Pole prostokąta jest równe  \( 1\,269\,\mathrm{cm}^2 \). Jego długość jest o  \( 20\,\mathrm{cm} \) większa niż szerokość. Wyznacz sumę długości i szerokości tego prostokąta.}
{\( 74\,\mathrm{cm} \)}{\( 35\,\mathrm{cm} \)}{\( 27\,\mathrm{cm} \)}{\( 57\,\mathrm{cm} \)}{}{}}
\LANGcs{
\Question{
Obdélník má obsah \( 1\,269\,\mathrm{cm}^2 \). Jeho délka je o \( 20\,\mathrm{cm} \) větší než šířka. Určete součet délky a šířky obdélníku.}
{\( 74\,\mathrm{cm} \)}{\( 35\,\mathrm{cm} \)}{\( 27\,\mathrm{cm} \)}{\( 57\,\mathrm{cm} \)}{}{}}
\LANGsk{
\Question{
Obdĺžnik má obsah \( 1\,269\,\mathrm{cm}^2 \). Jeho dĺžka je o \( 20\,\mathrm{cm} \) väčšia než jeho šírka. Určte súčet dĺžky a šírky daného obdĺžnika.}
{\( 74\,\mathrm{cm} \)}{\( 35\,\mathrm{cm} \)}{\( 27\,\mathrm{cm} \)}{\( 57\,\mathrm{cm} \)}{}{}}
\LANGes{
\Question{
El área de un rectángulo es \( 1\,269\,\mathrm{cm}^2 \). La longitud de un lado es \( 20\,\mathrm{cm} \) más larga que la longitud de otro lado. Halla la suma de los dos lados.}
{\( 74\,\mathrm{cm} \)}{\( 35\,\mathrm{cm} \)}{\( 27\,\mathrm{cm} \)}{\( 57\,\mathrm{cm} \)}{}{}}
\End

\Data\ID{1003085410}
\Updated{2020-08-22 04:08:13}
\Level{B}
\SubArea{Quadratic equations and inequalities}
\LANGen{
\Question{
Veronica and Joseph are on their romantic hot air balloon ride. When proposing in \( 1\,500\,\mathrm{m} \) altitude, Joseph accidentally drops a ring over the edge of the balloon basket. The ring perishes in the ground direction the same as does Veronicas  enthusiasm for getting married.   In how many seconds does the ring hit the ground? (Ignore air resistance and round the result to the nearest integer.) 

(Note: Distance \( d \) (in meters) travelled by an object falling for time \( t \) (in seconds) is given by  \( d=\frac12\,\mathrm{gt}^2 \), where \( g \) is gravitational acceleration, \( g = 9.81\,\mathrm{m/s^2} \).)}
{\( 17 \)}{\( 15 \)}{\( 20 \)}{\( 21 \)}{}{}}
\LANGpl{
\Question{
Weronika i Józef odbywają romantyczny lot balonem. Podczas oświadczyn na wysokości \( 1\,500\,\mathrm{m} \), Józef upuszcza pierścionek poza brzeg kosza. Pierścionek spada w kierunku ziemi tak samo jak entuzjazm Weroniki do zamążpójścia.   Za ile sekund pierścionek dotknie ziemi? (Nie bierz pod uwagę oporu powietrza i zaokrąglij wynik do najbliższej liczby całkowitej.) 

(Wskazówka: Odległość \( d \) (w metrach), którą przebył spadający przedmiot w czasie \( t \) (w sekundach) wyrażono przez  \( d=\frac12\,\mathrm{gt}^2 \), gdzie \( g \) jest przyspieszeniem ziemskim, \( g = 9{,}81\,\mathrm{m/s^2} \).)}
{\( 17 \)}{\( 15 \)}{\( 20 \)}{\( 21 \)}{}{}}
\LANGcs{
\Question{
Veronika s Josefem si zaplatili romantický let v horkovzdušném balónu.  Při žádosti o ruku ve výšce \( 1\,500\,\mathrm{m} \) Josef omylem upustil prsten přes okraj balónu a ten se řítí k zemi stejně rychle jako naděje Veroniky na sňatek. Za kolik sekund dopadne prsten na zem? Výsledek zaokrouhlete na celé číslo, odpor vzduchu zanedbejte.

(Poznámka: Volný pád je rovnoměrně zrychlený přímočarý pohyb se zrychlením rovným tíhovému zrychlení \( g = 9{,}81\,\mathrm{m/s^2} \). Dráhu \( s \) v čase \( t \) tělesa pohybujícího se volným pádem spočítáme dle  \( s=\frac12\,\mathrm{gt}^2 \).)}
{\( 17 \)}{\( 15 \)}{\( 20 \)}{\( 21 \)}{}{}}
\LANGsk{
\Question{
Veronika s Josefom podnikli romantický let v teplovzdušnom balóne. Pri žiadosti o ruku vo výške \( 1\,500\,\mathrm{m} \) Josef omylom pustil prsteň cez okraj balónu a ten sa rúti k zemi rovnako rýchle ako Veronikine nádeje na sobáš. Za koľko sekúnd dopadne prsteň na zem? (Odpor vzduchu zanedbajte a výsledok zaokrúhlite na  celé číslo.) 

(Poznámka: Dráhu \( s \) (v metroch) telesa pohybujúceho sa voľným pádom v čase \( t \) (v sekundách) vypočítame podľa  \( s=\frac12\,\mathrm{gt}^2 \), kde \( g \) is gravitačné zrýchlenie \( g = 9{,}81\,\mathrm{m/s^2} \).)}
{\( 17 \)}{\( 15 \)}{\( 20 \)}{\( 21 \)}{}{}}
\LANGes{
\Question{
Verónica y José están en un vuelo romántico en globo. A la altura de \( 1\,500\,\mathrm{m} \) José le pide la mano a Verónica pero desgraciadamente el anillo se le cae del globo. El anillo se despeña al suelo tan rápido como las ilusiones de Verónica de casarse. ¿Cuántos segundos tarda el anillo en caerse al suelo? Ignora la resistencia aerodinámica y aproxima el resultado al número entero más cercano. (Nota: Caída libre es  un caso particular del movimiento rectilíneo uniformemente acelerado igual a la intensidad del campo gravitatorio \( g = 9{,}81\,\mathrm{m/s^2} \). La trayectoria \( s \) en el tiempo \( t \) de un objeto en caída libre se calcula con la fórmula  \( s=\frac12\,\mathrm{gt}^2 \).)}
{\( 17 \)}{\( 15 \)}{\( 20 \)}{\( 21 \)}{}{}}
\End

\Data\ID{1003124101}
\Updated{2020-08-22 04:08:15}
\Level{B}
\SubArea{Quadratic equations and inequalities}
\LANGen{
\Question{
The following equation has roots \( -4 \) and \( 2 \). Find the value of the linear coefficient \( b \).
\[ x^2+bx-8=0 \]}
{\( 2 \)}{\( 4 \)}{\( 1 \)}{\( -2 \)}{}{}}
\LANGpl{
\Question{
Podane równanie ma pierwiastki  \( -4 \) i \( 2 \).  Wyznacz wartość liniowego współczynnika \( b \).
\[ x^2+bx-8=0 \]}
{\( 2 \)}{\( 4 \)}{\( 1 \)}{\( -2 \)}{}{}}
\LANGcs{
\Question{
Následující rovnice má kořeny  \( -4 \) a \( 2 \). Určete koeficient \( b \) lineárního členu dané rovnice.
\[ x^2+bx-8=0 \]}
{\( 2 \)}{\( 4 \)}{\( 1 \)}{\( -2 \)}{}{}}
\LANGsk{
\Question{
Nasledujúca rovnica má korene \( -4 \) a \( 2 \). Určte koeficient \( b \) lineárneho člena danej rovnice.
\[ x^2+bx-8=0 \]}
{\( 2 \)}{\( 4 \)}{\( 1 \)}{\( -2 \)}{}{}}
\LANGes{
\Question{
La siguiente ecuación tiene como raíces \( -4 \) y \( 2 \). Encuentra el valor del coeficiente \( b \) del miembro lineal de la ecuación 
\[ x^2+bx-8=0 \].}
{\( 2 \)}{\( 4 \)}{\( 1 \)}{\( -2 \)}{}{}}
\End

\Data\ID{1003124102}
\Updated{2020-08-22 04:08:37}
\Level{B}
\SubArea{Quadratic equations and inequalities}
\LANGen{
\Question{
The following equation has roots \( -\sqrt3 \) and \( \sqrt3 \). Find the value of the linear coefficient \( b \).
\[ x^2+bx+c=0 \]}
{\( 0 \)}{\( 1 \)}{\( 3 \)}{\( -3 \)}{}{}}
\LANGpl{
\Question{
Podane równanie ma pierwiastki \( -\sqrt3 \) i \( \sqrt3 \). Wyznacz wartość liniowego współczynnika \( b \).
\[ x^2+bx+c=0 \]}
{\( 0 \)}{\( 1 \)}{\( 3 \)}{\( -3 \)}{}{}}
\LANGcs{
\Question{
Následující rovnice má kořeny  \( -\sqrt3 \) a \( \sqrt3 \). Určete koeficient  \( b \) lineárního členu dané rovnice.
\[ x^2+bx+c=0 \]}
{\( 0 \)}{\( 1 \)}{\( 3 \)}{\( -3 \)}{}{}}
\LANGsk{
\Question{
Nasledujúca rovnica má korene \( -\sqrt3 \) a \( \sqrt3 \). Určte koeficient \( b \) lineárneho člena danej rovnice.
\[ x^2+bx+c=0 \]}
{\( 0 \)}{\( 1 \)}{\( 3 \)}{\( -3 \)}{}{}}
\LANGes{
\Question{
La siguiente ecuación tiene como raíces \( -\sqrt3 \) y \( \sqrt3 \). Encuentra el valor del coeficiente lineal \( b \).
\[ x^2+bx+c=0 \]}
{\( 0 \)}{\( 1 \)}{\( 3 \)}{\( -3 \)}{}{}}
\End

\Data\ID{1003124103}
\Updated{2020-08-22 04:08:06}
\Level{B}
\SubArea{Quadratic equations and inequalities}
\LANGen{
\Question{
The following equation has roots \( -1 \) and \( 5 \).  Find the value of the constant \( c \).
\[ x^2-4x+c=0 \]}
{\( -5 \)}{\( 4 \)}{\( -1 \)}{\( 1 \)}{}{}}
\LANGpl{
\Question{
Podane równanie zawiera pierwiastki  \( -1 \) i \( 5 \).  Wyznacz wartość stałej \( c \).
\[ x^2-4x+c=0 \]}
{\( -5 \)}{\( 4 \)}{\( -1 \)}{\( 1 \)}{}{}}
\LANGcs{
\Question{
Následující rovnice má kořeny  \( -1 \) a \( 5 \). Určete absolutní člen \( c \) dané rovnice.
\[ x^2-4x+c=0 \]}
{\( -5 \)}{\( 4 \)}{\( -1 \)}{\( 1 \)}{}{}}
\LANGsk{
\Question{
Nasledujúca rovnica má korene \( -1 \) a \( 5 \).  Určte absolútny člen \( c \) danej rovnice.
\[ x^2-4x+c=0 \]}
{\( -5 \)}{\( 4 \)}{\( -1 \)}{\( 1 \)}{}{}}
\LANGes{
\Question{
La siguiente ecuación tiene como raíces \( -1 \) y \( 5 \).  Halla el valor de la constante \( c \).
\[ x^2-4x+c=0 \]}
{\( -5 \)}{\( 4 \)}{\( -1 \)}{\( 1 \)}{}{}}
\End

\Data\ID{1003124104}
\Updated{2020-08-22 04:08:26}
\Level{B}
\SubArea{Quadratic equations and inequalities}
\LANGen{
\Question{
The following equation has roots \( -3 \) and \( -2 \). Find the value of the constant \( c \).
\[ x^2+bx+c=0 \]}
{\( 6 \)}{\( -3 \)}{\( -2 \)}{\( -6 \)}{}{}}
\LANGpl{
\Question{
Podane równanie posiada pierwiastki \( -3 \) i \( -2 \). Wyznacz wartość stałej \( c \).
\[ x^2+bx+c=0 \]}
{\( 6 \)}{\( -3 \)}{\( -2 \)}{\( -6 \)}{}{}}
\LANGcs{
\Question{
Následující rovnice má kořeny  \( -3 \) a \( -2 \). Určete absolutní člen \( c \) dané rovnice.
\[ x^2+bx+c=0 \]}
{\( 6 \)}{\( -3 \)}{\( -2 \)}{\( -6 \)}{}{}}
\LANGsk{
\Question{
Nasledujúca rovnica má korene \( -3 \) a \( -2 \). Určte absolútny člen \( c \) danej rovnice.
\[ x^2+bx+c=0 \]}
{\( 6 \)}{\( -3 \)}{\( -2 \)}{\( -6 \)}{}{}}
\LANGes{
\Question{
La siguiente ecuación tiene como raíces \( -3 \) y \( -2 \). Encuentra el valor de la constante \( c \).
\[ x^2+bx+c=0 \]}
{\( 6 \)}{\( -3 \)}{\( -2 \)}{\( -6 \)}{}{}}
\End

\Data\ID{1003124105}
\Updated{2020-08-09 08:08:16}
\Level{B}
\SubArea{Quadratic equations and inequalities}
\LANGen{
\Question{
Find the sum of the roots of the following equation. Use Vieta’s formulas, don’t solve the equation.
\[ x^2-x-12=0 \]}
{\( 1 \)}{\( -1 \)}{\( -12 \)}{\( 12 \)}{}{}}
\LANGpl{
\Question{
Wyznacz sumę wszystkich pierwiastków podanego równania. Użyj wzorów Vieta, ale nie rozwiązuj równania.
\[ x^2-x-12=0 \]}
{\( 1 \)}{\( -1 \)}{\( -12 \)}{\( 12 \)}{}{}}
\LANGcs{
\Question{
Určete součet kořenů následující rovnice. Zkuste příklad vyřešit s využitím Viètových vzorců, bez výpočtů kořenů rovnice.
\[ x^2-x-12=0 \]}
{\( 1 \)}{\( -1 \)}{\( -12 \)}{\( 12 \)}{}{}}
\LANGsk{
\Question{
Určte súčet koreňov nasledujúcej rovnice. Skúste príklad vyriešiť s využitím Vietových vzorcov, bez výpočtu koreňov rovnice.
\[ x^2-x-12=0 \]}
{\( 1 \)}{\( -1 \)}{\( -12 \)}{\( 12 \)}{}{}}
\LANGes{
\Question{
Halla la suma de las raíces de la siguiente ecuación. Usa las fórmulas de Cardano-Vieta sin calcular las raíces. 
\[ x^2-x-12=0 \]}
{\( 1 \)}{\( -1 \)}{\( -12 \)}{\( 12 \)}{}{}}
\End

\Data\ID{1003124106}
\Updated{2020-08-20 11:08:33}
\Level{B}
\SubArea{Quadratic equations and inequalities}
\LANGen{
\Question{
Find the roots of the following equation using Vieta’s formulas.
\[ x^2-8x+15=0 \]}
{\( \{3;5\} \)}{\( \{ -3;-5\} \)}{\( \{-5;3 \} \)}{\( \{-3;5 \} \)}{}{}}
\LANGpl{
\Question{
Wyznacz pierwiastki podanego równania używając wzorów Vieta.
\[ x^2-8x+15=0 \]}
{\( \{3;5\} \)}{\( \{ -3;-5\} \)}{\( \{-5;3 \} \)}{\( \{-3;5 \} \)}{}{}}
\LANGcs{
\Question{
Zpaměti, za pomoci Viètových vzorců, určete množinu kořenů následující rovnice.
\[ x^2-8x+15=0 \]}
{\( \{3;5\} \)}{\( \{ -3;-5\} \)}{\( \{-5;3 \} \)}{\( \{-3;5 \} \)}{}{}}
\LANGsk{
\Question{
Určte množinu koreňov nasledujúcej rovnice s využitím Vietových vzorcov.
\[ x^2-8x+15=0 \]}
{\( \{3;5\} \)}{\( \{ -3;-5\} \)}{\( \{-5;3 \} \)}{\( \{-3;5 \} \)}{}{}}
\LANGes{
\Question{
Halla las raíces de la siguiente ecuación utilizando las fórmulas de Cardano-Vieta.
\[ x^2-8x+15=0 \]}
{\( \{3;5\} \)}{\( \{ -3;-5\} \)}{\( \{-5;3 \} \)}{\( \{-3;5 \} \)}{}{}}
\End

\Data\ID{1003124107}
\Updated{2020-08-09 07:08:06}
\Level{B}
\SubArea{Quadratic equations and inequalities}
\LANGen{
\Question{
Find the roots of the following equation using Vieta’s formulas.
\[ x^2-11x+10=0 \]}
{\( \{1; 10 \} \)}{\( \{ -11; -1\} \)}{\( \{1;11 \} \)}{\( \{-1; 10 \} \)}{}{}}
\LANGpl{
\Question{
Wyznacz pierwiastki podanego równania używając wzorów Vieta.
\[ x^2-11x+10=0 \]}
{\( \{1; 10 \} \)}{\( \{ -11; -1\} \)}{\( \{1;11 \} \)}{\( \{-1; 10 \} \)}{}{}}
\LANGcs{
\Question{
Zpaměti, za pomoci Viètových vzorců, určete množinu kořenů následující rovnice.
\[ x^2-11x+10=0 \]}
{\( \{1; 10 \} \)}{\( \{ -11; -1\} \)}{\( \{1;11 \} \)}{\( \{-1; 10 \} \)}{}{}}
\LANGsk{
\Question{
Určte množinu koreňov nasledujúcej rovnice s využitím Vietových vzorcov.
\[ x^2-11x+10=0 \]}
{\( \{1; 10 \} \)}{\( \{ -11; -1\} \)}{\( \{1;11 \} \)}{\( \{-1; 10 \} \)}{}{}}
\LANGes{
\Question{
Encuentra las raíces de la siguiente ecuación usando las fórmulas de Cardano-Vieta. 
\[ x^2-11x+10=0 \]}
{\( \{1; 10 \} \)}{\( \{ -11; -1\} \)}{\( \{1;11 \} \)}{\( \{-1; 10 \} \)}{}{}}
\End

\Data\ID{9000020406}
\Updated{2020-08-22 04:08:22}
\Level{B}
\SubArea{Quadratic equations and inequalities}
\LANGen{
\Question{
The ratio of the sides of a rectangle is
\(3 : 4\). The length of
the diagonal is \(100\, \mathrm{cm}\).
Find the perimeter of the rectangle.}
{\(280\, \mathrm{cm}\)}{\(150\, \mathrm{cm}\)}{\(480\, \mathrm{cm}\)}{\(300\, \mathrm{cm}\)}{}{}}
\LANGpl{
\Question{
Stosunek boków prostokąta jest równy
\(3 : 4\). Długość przekątnej wynosi \(100\, \mathrm{cm}\).
Jaki jest obwód tego prostokąta?}
{\(280\, \mathrm{cm}\)}{\(150\, \mathrm{cm}\)}{\(480\, \mathrm{cm}\)}{\(300\, \mathrm{cm}\)}{}{}}
\LANGcs{
\Question{
Poměr délek stran obdélníku je \(3 : 4\).
Úhlopříčka měří \(100\, \mathrm{cm}\).
Určete obvod obdélníku v centimetrech.}
{\(280\, \mathrm{cm}\)}{\(150\, \mathrm{cm}\)}{\(480\, \mathrm{cm}\)}{\(300\, \mathrm{cm}\)}{}{}}
\LANGsk{
\Question{
Pomer dĺžok strán obdĺžnika je
\(3 : 4\). Uhlopriečka meria
 \(100\, \mathrm{cm}\).
Určte obvod obdĺžnika v cm.}
{\(280\, \mathrm{cm}\)}{\(150\, \mathrm{cm}\)}{\(480\, \mathrm{cm}\)}{\(300\, \mathrm{cm}\)}{}{}}
\LANGes{
\Question{
La proporción de los lados de un rectángulo es 
\(3 : 4\). La longitud de la diagonal es \(100\, \mathrm{cm}\).
Halla el perímetro del rectángulo.}
{\(280\, \mathrm{cm}\)}{\(150\, \mathrm{cm}\)}{\(480\, \mathrm{cm}\)}{\(300\, \mathrm{cm}\)}{}{}}
\End

\Data\ID{9000020409}
\Updated{2021-08-19 03:08:39}
\Level{B}
\SubArea{Quadratic equations and inequalities}
\LANGen{
\Question{
One of the solutions of the quadratic equation \( x^{2} + bx - 10 = 0\) is \(x_{1} = 5\). Find the second solution \(x_{2}\)  and the value of the coefficient \(b\).}
{\(x_{2} = -2\) and
\(b = -3\)}{\(x_{2} = -3\) and
\(b = -2\)}{\(x_{2} = 2\) and
\(b = 3\)}{\(x_{2} = 3\) and
\(b = 2\)}{}{}}
\LANGpl{
\Question{
Równanie kwadratowe 
\[ 
 x^{2} + bx - 10 = 0
\]
ma jedno rozwiązanie  \(x_{1} = 5\). Znajdź drugie rozwiązanie \(x_{2}\) 
i współczynnik  \(b\).}
{\(x_{2} = -2\) i \(b = -3\)}{\(x_{2} = -3\) i \(b = -2\)}{\(x_{2} = 2\) i \(b = 3\)}{\(x_{2} = 3\) i \(b = 2\)}{}{}}
\LANGcs{
\Question{
Jeden kořen kvadratické rovnice \[x^{2} + bx - 10 = 0\] 
je \(x_{1} = 5\).
Určete hodnotu druhého kořenu a hodnotu koeficientu
\(b\).}
{\(x_{2} = -2\) a
\(b = -3\)}{\(x_{2} = -3\) a
\(b = -2\)}{\(x_{2} = 2\) a
\(b = 3\)}{\(x_{2} = 3\) a
\(b = 2\)}{}{}}
\LANGsk{
\Question{
Jeden koreň kvadratickej rovnice  
\[ 
 x^{2} + bx - 10 = 0
\]
je \(x_{1} = 5\). Určte
hodnotu druhého koreňa \(x_{2}\) 
a hodnotu koeficientu \(b\).}
{\(x_{2} = -2\) a
\(b = -3\)}{\(x_{2} = -3\) a
\(b = -2\)}{\(x_{2} = 2\) a
\(b = 3\)}{\(x_{2} = 3\) a
\(b = 2\)}{}{}}
\LANGes{
\Question{
La ecuación cuadrática  
\[ 
 x^{2} + bx - 10 = 0
\]
tiene una solución \(x_{1} = 5\). Halla 
la segunda solución \(x_{2}\) 
y el coeficiente \(b\).}
{\(x_{2} = -2\) y
\(b = -3\)}{\(x_{2} = -3\) y
\(b = -2\)}{\(x_{2} = 2\) y
\(b = 3\)}{\(x_{2} = 3\) y
\(b = 2\)}{}{}}
\End

\Data\ID{9000020410}
\Updated{2020-08-20 08:08:02}
\Level{B}
\SubArea{Quadratic equations and inequalities}
\LANGen{
\Question{
The quadratic equation 
\[ 
 ax^{2} + 4x + c = 0
\]
has solutions \(x_{1} = -3\) 
and \(x_{2} = 5\). Find the
coefficients \(a\) 
and \(c\).}
{\(a = -2\),
\(c = 30\)}{\(a = -2\),
\(c = -30\)}{\(a = 2\),
\(c = -30\)}{\(a = 2\),
\(c = 30\)}{}{}}
\LANGpl{
\Question{
Równanie kwadratowe
\[ 
 ax^{2} + 4x + c = 0
\]
ma następujące rozwiązania \(x_{1} = -3\) 
i \(x_{2} = 5\). Znajdź współczynniki \(a\) 
i \(c\).}
{\(a = -2\),
\(c = 30\)}{\(a = -2\),
\(c = -30\)}{\(a = 2\),
\(c = -30\)}{\(a = 2\),
\(c = 30\)}{}{}}
\LANGcs{
\Question{
Kvadratická rovnice \[ax^{2} + 4x + c = 0\]
má kořeny \(- 3\) a
\(5\). Určete hodnotu
koeficientů \(a\),
\(c\).}
{\(a = -2,\ c = 30\)}{\(a = -2,\ c = -30\)}{\(a = 2,\ c = -30\)}{\(a = 2,\ c = 30\)}{}{}}
\LANGsk{
\Question{
Kvadratická rovnica 
\[ 
 ax^{2} + 4x + c = 0
\]
má korene \(x_{1} = -3\) 
a \(x_{2} = 5\). Určte
hodnotu koeficientov \(a\) 
a \(c\).}
{\(a = -2\),
\(c = 30\)}{\(a = -2\),
\(c = -30\)}{\(a = 2\),
\(c = -30\)}{\(a = 2\),
\(c = 30\)}{}{}}
\LANGes{
\Question{
La ecuación cuadrática  
\[ 
 ax^{2} + 4x + c = 0
\]
tiene soluciones \(x_{1} = -3\) 
y \(x_{2} = 5\). Encuentra los
coeficientes \(a\) 
y \(c\).}
{\(a = -2\),
\(c = 30\)}{\(a = -2\),
\(c = -30\)}{\(a = 2\),
\(c = -30\)}{\(a = 2\),
\(c = 30\)}{}{}}
\End

\Data\ID{9000020909}
\Updated{2020-08-22 05:08:31}
\Level{B}
\SubArea{Quadratic equations and inequalities}
\LANGen{
\Question{
The sum of squares of two consecutive integers is
\(1201\).
Identify these integers.}
{\(24\) 
and \(25\)}{\(23\) 
and \(24\)}{\(25\) 
and \(26\)}{\(26\) 
and \(27\)}{}{}}
\LANGpl{
\Question{
Suma kwadratów dwóch kolejnych liczb całkowitych wynosi
\(1201\).
Znajdź te liczby całkowite.}
{\(24\) 
i \(25\)}{\(23\) 
i \(24\)}{\(25\) 
i \(26\)}{\(26\) 
i \(27\)}{}{}}
\LANGcs{
\Question{
Součet druhých mocnin dvou po sobě jdoucích přirozených čísel je
\(1201\).
Určete obě čísla.}
{\(24\) a
\(25\)}{\(23\) a
\(24\)}{\(25\) a
\(26\)}{\(26\) a
\(27\)}{}{}}
\LANGsk{
\Question{
Súčet druhých mocnín dvoch po sebe idúcich celých čísel je 
\(1201\).
Určte obe čísla.}
{\(24\) 
a \(25\)}{\(23\) 
a \(24\)}{\(25\) 
a \(26\)}{\(26\) 
a \(27\)}{}{}}
\LANGes{
\Question{
La suma de las segundas potencias de dos números naturales consecutivos es
\(1201\). Identifica estos dos números.}
{\(24\) 
y \(25\)}{\(23\) 
y \(24\)}{\(25\) 
y \(26\)}{\(26\) 
y \(27\)}{}{}}
\End

\Data\ID{9000021703}
\Updated{2020-08-20 08:08:39}
\Level{B}
\SubArea{Quadratic equations and inequalities}
\LANGen{
\Question{
Solve the following inequality. 
\[ 
 (x - 2)^{2}\geq (x + 1)(x - 5)
\]}
{\(x\in \mathbb{R}\)}{\(x\in \emptyset \)}{\(x\in \left (-\infty ; \frac{9} 
{8}\right ] \)}{\(x\in \left [ \frac{9}
{8};\infty \right )\)}{}{}}
\LANGpl{
\Question{
Rozwiąż podaną nierówność:
\[ 
 (x - 2)^{2}\geq (x + 1)(x - 5)
\]}
{\(x\in \mathbb{R}\)}{\(x\in \emptyset \)}{\(x\in \left (-\infty ; \frac{9} 
{8}\right ] \)}{\(x\in \left [ \frac{9}
{8};\infty \right )\)}{}{}}
\LANGcs{
\Question{
Vyřešte danou nerovnici.
\[(x - 2)^{2}\geq (x + 1)(x - 5)\]}
{\(x\in\mathbb{R}\)}{\(x\in\emptyset \)}{\(x\in\left (-\infty ; \frac{9} 
{8}\right \rangle \)}{\(x\in\left \langle \frac{9}
{8};\infty \right )\)}{}{}}
\LANGsk{
\Question{
Vyriešte danú nerovnicu.
\[ 
 (x - 2)^{2}\geq (x + 1)(x - 5)
\]}
{\(x\in \mathbb{R}\)}{\(x\in \emptyset \)}{\(x\in \left (-\infty ; \frac{9} 
{8}\right ] \)}{\(x\in \left [ \frac{9}
{8};\infty \right )\)}{}{}}
\LANGes{
\Question{
Resuelve la siguiente inecuación.  
\[ 
 (x - 2)^{2}\geq (x + 1)(x - 5)
\]}
{\(x\in \mathbb{R}\)}{\(x\in \emptyset \)}{\(x\in \left (-\infty ; \frac{9} 
{8}\right ] \)}{\(x\in \left [ \frac{9}
{8};\infty \right )\)}{}{}}
\End

\Data\ID{9000021803}
\Updated{2020-08-20 08:08:33}
\Level{B}
\SubArea{Quadratic equations and inequalities}
\LANGen{
\Question{
Solve the following inequality. 
\[ 
 (3x - 1)(2 - 4x) < 0
\]}
{\(x\in \left (-\infty ; \frac{1} 
{3}\right )\cup \left (\frac{1}
{2};\infty \right )\)}{\(x\in \left (\frac{1}
{3}; \frac{1} 
{2}\right )\)}{\(x\in \left (-\infty ; \frac{1} 
{2}\right )\)}{\(x\in \left (\frac{1}
{3};\infty \right )\)}{}{}}
\LANGpl{
\Question{
Rozwiąż podaną nierówność:
\[ 
 (3x - 1)(2 - 4x) < 0
\]}
{\(x\in \left (-\infty ; \frac{1} 
{3}\right )\cup \left (\frac{1}
{2};\infty \right )\)}{\(x\in \left (\frac{1}
{3}; \frac{1} 
{2}\right )\)}{\(x\in \left (-\infty ; \frac{1} 
{2}\right )\)}{\(x\in \left (\frac{1}
{3};\infty \right )\)}{}{}}
\LANGcs{
\Question{
Vyřešte danou nerovnici.
\[(3x - 1)(2 - 4x) < 0\]}
{\(x\in\left (-\infty ; \frac{1} 
{3}\right )\cup \left (\frac{1}
{2};\infty \right )\)}{\(x\in\left (\frac{1}
{3}; \frac{1} 
{2}\right )\)}{\(x\in\left (-\infty ; \frac{1} 
{2}\right )\)}{\(x\in\left (\frac{1}
{3};\infty \right )\)}{}{}}
\LANGsk{
\Question{
Vyriešte danú nerovnicu. 
\[ 
 (3x - 1)(2 - 4x) < 0
\]}
{\(x\in \left (-\infty ; \frac{1} 
{3}\right )\cup \left (\frac{1}
{2};\infty \right )\)}{\(x\in \left (\frac{1}
{3}; \frac{1} 
{2}\right )\)}{\(x\in \left (-\infty ; \frac{1} 
{2}\right )\)}{\(x\in \left (\frac{1}
{3};\infty \right )\)}{}{}}
\LANGes{
\Question{
Resuelve la siguiente inecuación.  
\[ 
 (3x - 1)(2 - 4x) < 0
\]}
{\(x\in \left (-\infty ; \frac{1} 
{3}\right )\cup \left (\frac{1}
{2};\infty \right )\)}{\(x\in \left (\frac{1}
{3}; \frac{1} 
{2}\right )\)}{\(x\in \left (-\infty ; \frac{1} 
{2}\right )\)}{\(x\in \left (\frac{1}
{3};\infty \right )\)}{}{}}
\End

\Data\ID{9000022301}
\Updated{2020-08-20 08:08:06}
\Level{B}
\SubArea{Quadratic equations and inequalities}
\LANGen{
\Question{
Find the solution set of the quadratic inequality. 
\[ 
 x^{2} - 8x + 16\leq 0
\]}
{\(\{4\}\)}{\(\emptyset \)}{\(\mathbb{R}\setminus \{4\}\)}{\(\mathbb{R}\)}{\((-\infty ;4)\cup (4;\infty )\)}{}}
\LANGpl{
\Question{
Znajdź zbiór rozwiązań podanej nierówności kwadratowej. 
\[ 
 x^{2} - 8x + 16\leq 0
\]}
{\(\{4\}\)}{\(\emptyset \)}{\(\mathbb{R}\setminus \{4\}\)}{\(\mathbb{R}\)}{\((-\infty ;4)\cup (4;\infty )\)}{}}
\LANGcs{
\Question{
Určete množinu všech řešení dané nerovnice.
\[x^{2} - 8x + 16\leq 0\]}
{\(\{4\}\)}{\(\emptyset \)}{\(\mathbb{R}\setminus \{4\}\)}{\(\mathbb{R}\)}{\((-\infty ;4)\cup (4;\infty )\)}{}}
\LANGsk{
\Question{
Určte množinu všetkých riešení dáne nerovnice. \[
 x^{2} - 8x + 16\leq 0
\]}
{\(\{4\}\)}{\(\emptyset \)}{\(\mathbb{R}\setminus \{4\}\)}{\(\mathbb{R}\)}{\((-\infty ;4)\cup (4;\infty )\)}{}}
\LANGes{
\Question{
Halla el conjunto de soluciones de la inecuación cuadrática.  
\[ 
 x^{2} - 8x + 16\leq 0
\]}
{\(\{4\}\)}{\(\emptyset \)}{\(\mathbb{R}\setminus \{4\}\)}{\(\mathbb{R}\)}{\((-\infty ;4)\cup (4;\infty )\)}{}}
\End

\Data\ID{9000022303}
\Updated{2020-08-21 05:08:58}
\Level{B}
\SubArea{Quadratic equations and inequalities}
\LANGen{
\Question{
The solution set of one of the following quadratic inequalities is the interval
\((2;5)\).
Determine this inequality.}
{\(x^{2} - 7x + 10 < 0\)}{\(x^{2} + 7x + 10 > 0\)}{\(x^{2} - 7x + 10\leq 0\)}{\(x^{2} + 7x + 10\geq 0\)}{\(x^{2} - 7x + 10 > 0\)}{}}
\LANGpl{
\Question{
Zbiór rozwiązań jednej z podanych nierówności kwadratowych mieści się w przedziale
\((2;5)\).
Która to nierówność?}
{\(x^{2} - 7x + 10 < 0\)}{\(x^{2} + 7x + 10 > 0\)}{\(x^{2} - 7x + 10\leq 0\)}{\(x^{2} + 7x + 10\geq 0\)}{\(x^{2} - 7x + 10 > 0\)}{}}
\LANGcs{
\Question{
Vyberte tu z nerovnic, jejíž množinou řešení je interval
\((2;5)\).}
{\(x^{2} - 7x + 10 < 0\)}{\(x^{2} + 7x + 10 > 0\)}{\(x^{2} - 7x + 10\leq 0\)}{\(x^{2} + 7x + 10\geq 0\)}{\(x^{2} - 7x + 10 > 0\)}{}}
\LANGsk{
\Question{
Vyberte tú z nerovníc, ktorej množinou riešení je interval
\((2;5)\).}
{\(x^{2} - 7x + 10 < 0\)}{\(x^{2} + 7x + 10 > 0\)}{\(x^{2} - 7x + 10\leq 0\)}{\(x^{2} + 7x + 10\geq 0\)}{\(x^{2} - 7x + 10 > 0\)}{}}
\LANGes{
\Question{
El conjunto de soluciones de una de las siguientes inecuaciones es el intervalo \((2;5)\).
Determina esta inecuación.}
{\(x^{2} - 7x + 10 < 0\)}{\(x^{2} + 7x + 10 > 0\)}{\(x^{2} - 7x + 10\leq 0\)}{\(x^{2} + 7x + 10\geq 0\)}{\(x^{2} - 7x + 10 > 0\)}{}}
\End

\Data\ID{9000022304}
\Updated{2020-08-22 05:08:23}
\Level{B}
\SubArea{Quadratic equations and inequalities}
\LANGen{
\Question{
Find all the values of \(x\) 
at which the following expression attains nonnegative value. 
\[ 
 x^{2} + x - 12
\]}
{\(x\in \left (-\infty ;-4\right ] \cup \left [ 3;\infty \right )\)}{\(x\in \left [ -3;4\right ] \)}{\(x\in \left [ -4;3\right ] \)}{\(x\in \left (-\infty ;-4\right )\cup \left (3;\infty \right )\)}{\(x\in \left (-\infty ;-3\right ] \cup \left [ 4;\infty \right )\)}{}}
\LANGpl{
\Question{
Znajdź wszystkie wartości  \(x\), dla których podane wyrażenie przyjmuje wartość nieujemną. 
\[ 
 x^{2} + x - 12
\]}
{\(x\in \left (-\infty ;-4\right ] \cup \left [ 3;\infty \right )\)}{\(x\in \left [ -3;4\right ] \)}{\(x\in \left [ -4;3\right ] \)}{\(x\in \left (-\infty ;-4\right )\cup \left (3;\infty \right )\)}{\(x\in \left (-\infty ;-3\right ] \cup \left [ 4;\infty \right )\)}{}}
\LANGcs{
\Question{
Vyberte všechna \(x\),
pro která je daný výraz
nezáporný.
\[x^{2} + x - 12\]}
{\(x\in \left (-\infty ;-4\right \rangle \cup \left \langle 3;\infty \right )\)}{\(x\in \left \langle -3;4\right \rangle \)}{\(x\in \left \langle -4;3\right \rangle \)}{\(x\in \left (-\infty ;-4\right )\cup \left (3;\infty \right )\)}{\(x\in \left (-\infty ;-3\right \rangle \cup \left \langle 4;\infty \right )\)}{}}
\LANGsk{
\Question{
Vyberte všetky hodnoty \(x\), pre ktoré je daný výraz nezáporný. 
\[ 
 x^{2} + x - 12
\]}
{\(x\in \left (-\infty ;-4\right ] \cup \left [ 3;\infty \right )\)}{\(x\in \left [ -3;4\right ] \)}{\(x\in \left [ -4;3\right ] \)}{\(x\in \left (-\infty ;-4\right )\cup \left (3;\infty \right )\)}{\(x\in \left (-\infty ;-3\right ] \cup \left [ 4;\infty \right )\)}{}}
\LANGes{
\Question{
Halla todos los valores de \(x\) 
para los que la siguiente expresión tiene valor no negativo.  
\[ 
 x^{2} + x - 12
\]}
{\(x\in \left (-\infty ;-4\right ] \cup \left [ 3;\infty \right )\)}{\(x\in \left [ -3;4\right ] \)}{\(x\in \left [ -4;3\right ] \)}{\(x\in \left (-\infty ;-4\right )\cup \left (3;\infty \right )\)}{\(x\in \left (-\infty ;-3\right ] \cup \left [ 4;\infty \right )\)}{}}
\End

\Data\ID{9000022805}
\Updated{2020-08-22 05:08:48}
\Level{B}
\SubArea{Quadratic equations and inequalities}
\LANGen{
\Question{
The solution set of one of the following inequalities is the interval
\([ 3;5] \).
Identify this inequality.}
{\(x^{2} - 8x + 15\leq 0\)}{\(x^{2} + 8x + 15\leq 0\)}{\(x^{2} - 8x + 15\geq 0\)}{\(x^{2} + 8x + 15\geq 0\)}{}{}}
\LANGpl{
\Question{
Zbiór rozwiązań jednej z podanych nierówności mieści się w przedziale 
\([ 3;5] \).
Która to nierówność?}
{\(x^{2} - 8x + 15\leq 0\)}{\(x^{2} + 8x + 15\leq 0\)}{\(x^{2} - 8x + 15\geq 0\)}{\(x^{2} + 8x + 15\geq 0\)}{}{}}
\LANGcs{
\Question{
Interval \(\langle 3;5\rangle \) 
je množinou všech řešení jedné z uvedených nerovnic. Určete tuto nerovnici.}
{\(x^{2} - 8x + 15\leq 0\)}{\(x^{2} + 8x + 15\leq 0\)}{\(x^{2} - 8x + 15\geq 0\)}{\(x^{2} + 8x + 15\geq 0\)}{}{}}
\LANGsk{
\Question{
Množina všetkých riešení jednej z nasledujúcich nerovníc je interval
\([ 3;5] \). Určte túto nerovnicu.}
{\(x^{2} - 8x + 15\leq 0\)}{\(x^{2} + 8x + 15\leq 0\)}{\(x^{2} - 8x + 15\geq 0\)}{\(x^{2} + 8x + 15\geq 0\)}{}{}}
\LANGes{
\Question{
El conjunto de soluciones de una de las siguientes inecuaciones es el intervalo 
\([ 3;5] \).
Identifica esta inecuación.}
{\(x^{2} - 8x + 15\leq 0\)}{\(x^{2} + 8x + 15\leq 0\)}{\(x^{2} - 8x + 15\geq 0\)}{\(x^{2} + 8x + 15\geq 0\)}{}{}}
\End

\Data\ID{9000022806}
\Updated{2020-08-21 05:08:21}
\Level{B}
\SubArea{Quadratic equations and inequalities}
\LANGen{
\Question{
For an integer variable \(x\) 
find the solution set of the following quadratic inequality. 
\[ 
 2x^{2} - x - 6\leq 0
\]}
{\(\{ - 1;0;1;2\}\)}{\(\{ - 2;-1;0;1\}\)}{\(\{0;1;2;3\}\)}{\(\{ - 3;-2;-1;0\}\)}{}{}}
\LANGpl{
\Question{
Dla zmiennej całkowitej  \(x\) 
znajdź zbiór rozwiązań następującej nierówności kwadratowej.
\[ 
 2x^{2} - x - 6\leq 0
\]}
{\(\{ - 1;0;1;2\}\)}{\(\{ - 2;-1;0;1\}\)}{\(\{0;1;2;3\}\)}{\(\{ - 3;-2;-1;0\}\)}{}{}}
\LANGcs{
\Question{
V oboru celých čísel najděte řešení dané kvadratické nerovnice.
\[2x^{2} - x - 6\leq 0\]}
{\(\{ - 1;0;1;2\}\)}{\(\{ - 2;-1;0;1\}\)}{\(\{0;1;2;3\}\)}{\(\{ - 3;-2;-1;0\}\)}{}{}}
\LANGsk{
\Question{
V obore celých čísel \(x\) 
nájdite množinu všetkých riešení danej kvadratickej nerovnice. 
\[ 
 2x^{2} - x - 6\leq 0
\]}
{\(\{ - 1;0;1;2\}\)}{\(\{ - 2;-1;0;1\}\)}{\(\{0;1;2;3\}\)}{\(\{ - 3;-2;-1;0\}\)}{}{}}
\LANGes{
\Question{
En el conjunto de números enteros encuentra el conjunto de soluciones de la siguiente inecuación cuadrática. 
\[ 
 2x^{2} - x - 6\leq 0
\]}
{\(\{ - 1;0;1;2\}\)}{\(\{ - 2;-1;0;1\}\)}{\(\{0;1;2;3\}\)}{\(\{ - 3;-2;-1;0\}\)}{}{}}
\End

\Data\ID{9000022807}
\Updated{2020-08-21 05:08:54}
\Level{B}
\SubArea{Quadratic equations and inequalities}
\LANGen{
\Question{
Complete the following statement: Quadratic inequality 
\[ 
 2x^{2} - 3x + 4 > x^{2} + 2x - 2
\]
is satisfied if and only if}
{\(x\in (-\infty ;2)\cup (3;\infty )\).}{\(x\in (2;3)\).}{\(x\in (-\infty ;-2)\cup (-3;\infty )\).}{\(x\in (-2;-3)\).}{}{}}
\LANGpl{
\Question{
Dokończ następujące stwierdzenie: Nierówność kwadratowa 
\[ 
 2x^{2} - 3x + 4 > x^{2} + 2x - 2
\]
jest spełniona tylko wtedy, gdy}
{\(x\in (-\infty ;2)\cup (3;\infty )\).}{\(x\in (2;3)\).}{\(x\in (-\infty ;-2)\cup (-3;\infty )\).}{\(x\in (-2;-3)\).}{}{}}
\LANGcs{
\Question{
Nerovnost \(2x^{2} - 3x + 4 > x^{2} + 2x - 2\) 
je splněna, právě když platí:}
{\(x\in (-\infty ;2)\cup (3;\infty )\)}{\(x\in (2;3)\)}{\(x\in (-\infty ;-2)\cup (-3;\infty )\)}{\(x\in (-2;-3)\)}{}{}}
\LANGsk{
\Question{
Kvadratická nerovnosť 
\[ 
 2x^{2} - 3x + 4 > x^{2} + 2x - 2
\]
je splnená, práve keď platí:}
{\(x\in (-\infty ;2)\cup (3;\infty )\)}{\(x\in (2;3)\)}{\(x\in (-\infty ;-2)\cup (-3;\infty )\)}{\(x\in (-2;-3)\)}{}{}}
\LANGes{
\Question{
Completa el siguiente enunciado: La inecuación cuadrática  
\[ 
 2x^{2} - 3x + 4 > x^{2} + 2x - 2
\]
se cumple solo si vale}
{\(x\in (-\infty ;2)\cup (3;\infty )\).}{\(x\in (2;3)\).}{\(x\in (-\infty ;-2)\cup (-3;\infty )\).}{\(x\in (-2;-3)\).}{}{}}
\End

\Data\ID{9000022808}
\Updated{2020-08-22 05:08:51}
\Level{B}
\SubArea{Quadratic equations and inequalities}
\LANGen{
\Question{
Find all the real values of \(x\) 
which ensure that the following expression is negative. 
\[ 
 -x^{2} + 4x - 4
\]}
{\(x\in \mathbb{R}\setminus \{2\}\)}{none \(x\) 
with this property}{\(x = 2\)}{\(x\in \mathbb{R}\)}{}{}}
\LANGpl{
\Question{
Znajdź wszystkie rzeczywiste wartości  \(x\), dla których następujące wyrażenie jest ujemne. 
\[ 
 -x^{2} + 4x - 4
\]}
{\(x\in \mathbb{R}\setminus \{2\}\)}{żadne  \(x\) 
z tymi własnościami}{\(x = 2\)}{\(x\in \mathbb{R}\)}{}{}}
\LANGcs{
\Question{
Pro která \(x\) 
je výraz \(- x^{2} + 4x - 4\) 
záporný?}
{pro všechna \(x\in \mathbb{R}\setminus \{2\}\)}{pro žádná \(x\)}{pro \(x = 2\)}{pro všechna \(x\in \mathbb{R}\)}{}{}}
\LANGsk{
\Question{
Nájdite všetky reálne hodnoty  \(x\), pre ktoré nasledujúci výraz je záporný.
 \[ 
 -x^{2} + 4x - 4
\]}
{\(x\in \mathbb{R}\setminus \{2\}\)}{žiadne \(x\) 
s touto vlastnosťou}{\(x = 2\)}{\(x\in \mathbb{R}\)}{}{}}
\LANGes{
\Question{
Encuentra todos los valores de \(x\) 
para los cuales se cumple que la siguiente expresión es negativa.  
\[ 
 -x^{2} + 4x - 4
\]}
{\(x\in \mathbb{R}\setminus \{2\}\)}{no existe \(x\) 
con estas propiedades}{\(x = 2\)}{\(x\in \mathbb{R}\)}{}{}}
\End

\Data\ID{9000022809}
\Updated{2020-08-20 08:08:04}
\Level{B}
\SubArea{Quadratic equations and inequalities}
\LANGen{
\Question{
Find the solution set of the following quadratic inequality. 
\[ 
 4x^{2} + 4x + 1 < 0
\]}
{\(\emptyset \)}{\(\mathbb{R}\)}{\(\left \{\frac{1}
{2}\right \}\)}{\(\left \{-\frac{1}
{2}\right \}\)}{}{}}
\LANGpl{
\Question{
Znajdź zbiór rozwiązań następującej nierówności kwadratowej.
\[ 
 4x^{2} + 4x + 1 < 0
\]}
{\(\emptyset \)}{\(\mathbb{R}\)}{\(\left \{\frac{1}
{2}\right \}\)}{\(\left \{-\frac{1}
{2}\right \}\)}{}{}}
\LANGcs{
\Question{
Množina všech řešení kvadratické nerovnice
\(4x^{2} + 4x + 1 < 0\) je:}
{\(\emptyset \)}{\(\mathbb{R}\)}{\(\left \{\frac{1}
{2}\right \}\)}{\(\left \{-\frac{1}
{2}\right \}\)}{}{}}
\LANGsk{
\Question{
Nájdite množinu všetkých riešení danej kvadratickej nerovnice. 
\[ 
 4x^{2} + 4x + 1 < 0
\]}
{\(\emptyset \)}{\(\mathbb{R}\)}{\(\left \{\frac{1}
{2}\right \}\)}{\(\left \{-\frac{1}
{2}\right \}\)}{}{}}
\LANGes{
\Question{
Encuentra el conjunto de soluciones de la siguiente inecuación cuadrática.  
\[ 
 4x^{2} + 4x + 1 < 0
\]}
{\(\emptyset \)}{\(\mathbb{R}\)}{\(\left \{\frac{1}
{2}\right \}\)}{\(\left \{-\frac{1}
{2}\right \}\)}{}{}}
\End

\Data\ID{9000022810}
\Updated{2020-08-20 08:08:34}
\Level{B}
\SubArea{Quadratic equations and inequalities}
\LANGen{
\Question{
Find the solution set of the following quadratic inequality. 
\[ 
 -x^{2} + 2x + 3 > 0
\]}
{\((-1;3)\)}{\((-\infty ;-1)\)}{\((-\infty ;-1)\cup (3;\infty )\)}{\((3;\infty )\)}{}{}}
\LANGpl{
\Question{
Znajdź zbiór rozwiązań podanej nierówności kwadratowej.
\[ 
 -x^{2} + 2x + 3 > 0
\]}
{\((-1;3)\)}{\((-\infty ;-1)\)}{\((-\infty ;-1)\cup (3;\infty )\)}{\((3;\infty )\)}{}{}}
\LANGcs{
\Question{
Množina všech řešení kvadratické nerovnice
\(- x^{2} + 2x + 3 > 0\) je:}
{\((-1;3)\)}{\((-\infty ;-1)\)}{\((-\infty ;-1)\cup (3;\infty )\)}{\((3;\infty )\)}{}{}}
\LANGsk{
\Question{
Nájdite množinu všetkých riešení danej kvadratickej nerovnice. 
\[ 
 -x^{2} + 2x + 3 > 0
\]}
{\((-1;3)\)}{\((-\infty ;-1)\)}{\((-\infty ;-1)\cup (3;\infty )\)}{\((3;\infty )\)}{}{}}
\LANGes{
\Question{
Encuentra el conjunto de soluciones de la siguiente inecuación cuadrática.  
\[ 
 -x^{2} + 2x + 3 > 0
\]}
{\((-1;3)\)}{\((-\infty ;-1)\)}{\((-\infty ;-1)\cup (3;\infty )\)}{\((3;\infty )\)}{}{}}
\End

\Data\ID{9000033701}
\Updated{2021-08-19 03:08:55}
\Level{B}
\SubArea{Quadratic equations and inequalities}
\LANGen{
\Question{
Find out, how many integer solutions the following inequality has.
\[ 
 m^{2} + 2m - 4 < 0
\]}
{Five integer solutions.}{Less than five integer solutions.}{More than five integer solutions.}{}{}{}}
\LANGpl{
\Question{
Ile liczb całkowitych spełnia podaną nierówność? 
\[ 
 m^{2} + 2m - 4 < 0
\]}
{Pięć liczb całkowitych.}{Mniej niż pięć liczb całkowitych.}{Więcej niż pięć liczb całkowitych.}{}{}{}}
\LANGcs{
\Question{
Rozhodněte o počtu celočíselných řešení následující nerovnice.
\[m^{2} + 2m - 4 < 0\]}
{Nerovnice má právě pět řešení.}{Nerovnice má méně než pět řešení.}{Nerovnice má více než pět řešení.}{}{}{}}
\LANGsk{
\Question{
Rozhodnite o počte riešení danej nerovnice, kde koreň je celé číslo.
\[ 
 m^{2} + 2m - 4 < 0
\]}
{Päť celočíselných riešení.}{Menej ako päť celočíselných riešení.}{Viac ako päť celočíselných riešení.}{}{}{}}
\LANGes{
\Question{
Encuentra el número de soluciones de la siguiente inecuación. 
\[ 
 m^{2} + 2m - 4 < 0
\]}
{La inecuación tiene cinco soluciones.}{La inecuación tiene menos de cinco soluciones.}{La inecuación tiene más de cinco soluciones.}{}{}{}}
\End

\Data\ID{9000034905}
\Updated{2020-08-21 07:08:34}
\Level{B}
\SubArea{Quadratic equations and inequalities}
\LANGen{
\Question{
The solution set of one of the following quadratic inequalities is the interval
\(\left [ -\frac{7}
{6}; \frac{3} 
{4}\right ] \).
Determine this inequality.}
{\(\left (x + \frac{7} 
{6}\right )\left (x -\frac{3} 
{4}\right )\leq 0\)}{\(\left (x + \frac{7} 
{6}\right )\left (x -\frac{3} 
{4}\right )\geq 0\)}{\(\left (x -\frac{7} 
{6}\right )\left (x + \frac{3} 
{4}\right )\geq 0\)}{\(\left (x -\frac{7} 
{6}\right )\left (x + \frac{3} 
{4}\right )\leq 0\)}{}{}}
\LANGpl{
\Question{
Zbiór rozwiązań jednej z następujących nierówności kwadratowych mieści się w przedziale  \(\left [ -\frac{7}
{6}; \frac{3} 
{4}\right ] \).
Która to nierówność?}
{\(\left (x + \frac{7} 
{6}\right )\left (x -\frac{3} 
{4}\right )\leq 0\)}{\(\left (x + \frac{7} 
{6}\right )\left (x -\frac{3} 
{4}\right )\geq 0\)}{\(\left (x -\frac{7} 
{6}\right )\left (x + \frac{3} 
{4}\right )\geq 0\)}{\(\left (x -\frac{7} 
{6}\right )\left (x + \frac{3} 
{4}\right )\leq 0\)}{}{}}
\LANGcs{
\Question{
Určete kvadratickou nerovnici, jejíž množinou řešení je interval \(\left \langle -\frac{7}
{6}; \frac{3} 
{4}\right \rangle \).}
{\(\left (x + \frac{7} 
{6}\right )\left (x -\frac{3} 
{4}\right )\leq 0\)}{\(\left (x + \frac{7} 
{6}\right )\left (x -\frac{3} 
{4}\right )\geq 0\)}{\(\left (x -\frac{7} 
{6}\right )\left (x + \frac{3} 
{4}\right )\geq 0\)}{\(\left (x -\frac{7} 
{6}\right )\left (x + \frac{3} 
{4}\right )\leq 0\)}{}{}}
\LANGsk{
\Question{
Množina všetkých riešení kvadratickej nerovnice je interval
\(\left [ -\frac{7}
{6}; \frac{3} 
{4}\right ] \).
Určte túto nerovnicu.}
{\(\left (x + \frac{7} 
{6}\right )\left (x -\frac{3} 
{4}\right )\leq 0\)}{\(\left (x + \frac{7} 
{6}\right )\left (x -\frac{3} 
{4}\right )\geq 0\)}{\(\left (x -\frac{7} 
{6}\right )\left (x + \frac{3} 
{4}\right )\geq 0\)}{\(\left (x -\frac{7} 
{6}\right )\left (x + \frac{3} 
{4}\right )\leq 0\)}{}{}}
\LANGes{
\Question{
Determina la inecuación cuadrática cuyo conjunto de soluciones es el intervalo \(\left [ -\frac{7}
{6}; \frac{3} 
{4}\right ] \).}
{\(\left (x + \frac{7} 
{6}\right )\left (x -\frac{3} 
{4}\right )\leq 0\)}{\(\left (x + \frac{7} 
{6}\right )\left (x -\frac{3} 
{4}\right )\geq 0\)}{\(\left (x -\frac{7} 
{6}\right )\left (x + \frac{3} 
{4}\right )\geq 0\)}{\(\left (x -\frac{7} 
{6}\right )\left (x + \frac{3} 
{4}\right )\leq 0\)}{}{}}
\End

\Data\ID{9000034906}
\Updated{2020-08-21 07:08:11}
\Level{B}
\SubArea{Quadratic equations and inequalities}
\LANGen{
\Question{
The solution set of one of the following quadratic inequalities is
\(\left (-\infty ;-\frac{3}
{5}\right )\cup \left (\frac{1}
{6};\infty \right )\). Determine this inequality.}
{\(\left (5x + 3\right )\left (1 - 6x\right ) < 0\)}{\(\left (5x - 3\right )\left (6x + 1\right ) < 0\)}{\(\left (5x + 3\right )\left (1 - 6x\right ) > 0\)}{\(\left (5x - 3\right )\left (6x + 1\right ) > 0\)}{}{}}
\LANGpl{
\Question{
Zbiorem rozwiązań jednej z podanych nierówności jest przedział 
\(\left (-\infty ;-\frac{3}
{5}\right )\cup \left (\frac{1}
{6};\infty \right )\).
Znajdź tę nierówność.}
{\(\left (5x + 3\right )\left (1 - 6x\right ) < 0\)}{\(\left (5x - 3\right )\left (6x + 1\right ) < 0\)}{\(\left (5x + 3\right )\left (1 - 6x\right ) > 0\)}{\(\left (5x - 3\right )\left (6x + 1\right ) > 0\)}{}{}}
\LANGcs{
\Question{
Určete kvadratickou nerovnici, jejíž množinou řešení je interval \(\left (-\infty ;-\frac{3}
{5}\right )\cup \left (\frac{1}
{6};\infty \right )\).}
{\(\left (5x + 3\right )\left (1 - 6x\right ) < 0\)}{\(\left (5x - 3\right )\left (6x + 1\right ) < 0\)}{\(\left (5x + 3\right )\left (1 - 6x\right ) > 0\)}{\(\left (5x - 3\right )\left (6x + 1\right ) > 0\)}{}{}}
\LANGsk{
\Question{
Množina všetkých riešení kvadratickej nerovnice je 
\(\left (-\infty ;-\frac{3}
{5}\right )\cup \left (\frac{1}
{6};\infty \right )\).
Určte túto nerovnicu.}
{\(\left (5x + 3\right )\left (1 - 6x\right ) < 0\)}{\(\left (5x - 3\right )\left (6x + 1\right ) < 0\)}{\(\left (5x + 3\right )\left (1 - 6x\right ) > 0\)}{\(\left (5x - 3\right )\left (6x + 1\right ) > 0\)}{}{}}
\LANGes{
\Question{
Determina la inecuación cuadrática cuyo conjunto de soluciones es el intervalo
\(\left (-\infty ;-\frac{3}
{5}\right )\cup \left (\frac{1}
{6};\infty \right )\).}
{\(\left (5x + 3\right )\left (1 - 6x\right ) < 0\)}{\(\left (5x - 3\right )\left (6x + 1\right ) < 0\)}{\(\left (5x + 3\right )\left (1 - 6x\right ) > 0\)}{\(\left (5x - 3\right )\left (6x + 1\right ) > 0\)}{}{}}
\End

\Data\ID{9000034907}
\Updated{2020-08-21 07:08:35}
\Level{B}
\SubArea{Quadratic equations and inequalities}
\LANGen{
\Question{
Find all \(x\in \mathbb{R}\) 
for which the following expression takes nonnegative values. 
\[ 
 -2\left (x - 3\right )\left (2 - x\right )
\]}
{\(\left (-\infty ;2\right ] \cup \left [ 3;\infty \right )\)}{\(\left [ 2;3\right ] \)}{\(\left (2;3\right )\)}{\(\left (-\infty ;2\right )\cup \left (3;\infty \right )\)}{}{}}
\LANGpl{
\Question{
Znajdź wszystkie  \(x\in \mathbb{R}\),
dla których podane wyrażenie przyjmuje wartości nieujemne.
\[ 
 -2\left (x - 3\right )\left (2 - x\right )
\]}
{\(\left (-\infty ;2\right ] \cup \left [ 3;\infty \right )\)}{\(\left [ 2;3\right ] \)}{\(\left (2;3\right )\)}{\(\left (-\infty ;2\right )\cup \left (3;\infty \right )\)}{}{}}
\LANGcs{
\Question{
Množina všech \(x\in \mathbb{R}\), pro
která není výraz \(- 2\left (x - 3\right )\left (2 - x\right )\) 
záporný, je:}
{\(\left (-\infty ;2\right \rangle \cup \left \langle 3;\infty \right )\)}{\(\left \langle 2;3\right \rangle \)}{\(\left (2;3\right )\)}{\(\left (-\infty ;2\right )\cup \left (3;\infty \right )\)}{}{}}
\LANGsk{
\Question{
Nájdite všetky \(x\in \mathbb{R}\), pre ktoré daný výraz nadobúda nezáporné hodnoty.
 \[ 
 -2\left (x - 3\right )\left (2 - x\right )
\]}
{\(\left (-\infty ;2\right ] \cup \left [ 3;\infty \right )\)}{\(\left [ 2;3\right ] \)}{\(\left (2;3\right )\)}{\(\left (-\infty ;2\right )\cup \left (3;\infty \right )\)}{}{}}
\LANGes{
\Question{
Halla todos los \(x\in \mathbb{R}\) 
para los que la siguiente expresión no es negativa.  
\[ 
 -2\left (x - 3\right )\left (2 - x\right )
\]}
{\(\left (-\infty ;2\right ] \cup \left [ 3;\infty \right )\)}{\(\left [ 2;3\right ] \)}{\(\left (2;3\right )\)}{\(\left (-\infty ;2\right )\cup \left (3;\infty \right )\)}{}{}}
\End

\Data\ID{9000034908}
\Updated{2020-08-21 07:08:00}
\Level{B}
\SubArea{Quadratic equations and inequalities}
\LANGen{
\Question{
Find all \(x\in \mathbb{R}\) 
for which the following expression is nonpositive. 
\[ 
 \left (x + 1\right )\left (4 + x\right )
\]}
{\(\left [ -4;-1\right ] \)}{\(\left (-\infty ;-4\right ] \cup \left [ -1;\infty \right )\)}{\(\left (-4;-1\right )\)}{\(\left (-\infty ;-4\right )\cup \left (-1;\infty \right )\)}{}{}}
\LANGpl{
\Question{
Znajdź wszystkie \(x\in \mathbb{R}\), dla których podane wyrażenie jest niedodatnie. 
\[ 
 \left (x + 1\right )\left (4 + x\right )
\]}
{\(\left [ -4;-1\right ] \)}{\(\left (-\infty ;-4\right ] \cup \left [ -1;\infty \right )\)}{\(\left (-4;-1\right )\)}{\(\left (-\infty ;-4\right )\cup \left (-1;\infty \right )\)}{}{}}
\LANGcs{
\Question{
Množina všech \(x\in \mathbb{R}\), pro
která není výraz \(\left (x + 1\right )\left (4 + x\right )\) 
kladný, je:}
{\(\left \langle -4;-1\right \rangle \)}{\(\left (-\infty ;-4\right \rangle \cup \left \langle -1;\infty \right )\)}{\(\left (-4;-1\right )\)}{\(\left (-\infty ;-4\right )\cup \left (-1;\infty \right )\)}{}{}}
\LANGsk{
\Question{
Nájdite všetky \(x\in \mathbb{R}\), pre ktoré daný výraz nie je kladný. 
\[ 
 \left (x + 1\right )\left (4 + x\right )
\]}
{\(\left [ -4;-1\right ] \)}{\(\left (-\infty ;-4\right ] \cup \left [ -1;\infty \right )\)}{\(\left (-4;-1\right )\)}{\(\left (-\infty ;-4\right )\cup \left (-1;\infty \right )\)}{}{}}
\LANGes{
\Question{
Encuentra todos los \(x\in \mathbb{R}\) 
para los que la siguiente expresión no es positiva. 
\[ 
 \left (x + 1\right )\left (4 + x\right )
\]}
{\(\left [ -4;-1\right ] \)}{\(\left (-\infty ;-4\right ] \cup \left [ -1;\infty \right )\)}{\(\left (-4;-1\right )\)}{\(\left (-\infty ;-4\right )\cup \left (-1;\infty \right )\)}{}{}}
\End

\Data\ID{9000034909}
\Updated{2020-08-09 08:08:17}
\Level{B}
\SubArea{Quadratic equations and inequalities}
\LANGen{
\Question{
Find the solution set of the following quadratic inequality. 
\[ 
 -3(x + 2)^{2} < 0
\]}
{\(\mathbb{R}\setminus \{ - 2\}\)}{\(\emptyset \)}{\(\{- 2\}\)}{\(\mathbb{R}\)}{}{}}
\LANGpl{
\Question{
Znajdź zbiór rozwiązań następującej nierówności kwadratowej.
\[ 
 -3(x + 2)^{2} < 0
\]}
{\(\mathbb{R}\setminus \{ - 2\}\)}{\(\emptyset \)}{\(\{- 2\}\)}{\(\mathbb{R}\)}{}{}}
\LANGcs{
\Question{
Množina všech řešení kvadratické nerovnice
\(- 3(x + 2)^{2} < 0\) je:}
{\(\mathbb{R}\setminus \{ - 2\}\)}{\(\emptyset \)}{\(\{- 2\}\)}{\(\mathbb{R}\)}{}{}}
\LANGsk{
\Question{
Nájdite množinu všetkých riešení danej kvadratickej nerovnice. 
\[ 
 -3(x + 2)^{2} < 0
\]}
{\(\mathbb{R}\setminus \{ - 2\}\)}{\(\emptyset \)}{\(\{- 2\}\)}{\(\mathbb{R}\)}{}{}}
\LANGes{
\Question{
Encuentra el conjunto de soluciones de la siguiente inecuación cuadrática.  
\[ 
 -3(x + 2)^{2} < 0
\]}
{\(\mathbb{R}\setminus \{ - 2\}\)}{\(\emptyset \)}{\(\{- 2\}\)}{\(\mathbb{R}\)}{}{}}
\End

\Data\ID{9000034910}
\Updated{2020-08-20 08:08:13}
\Level{B}
\SubArea{Quadratic equations and inequalities}
\LANGen{
\Question{
Find the solution set of the following quadratic inequality. 
\[ 
 (x - 3)^{2}\geq 0
\]}
{\(\mathbb{R}\)}{\(\emptyset \)}{\(\mathbb{R}\setminus \{3\}\)}{\(\{3\}\)}{}{}}
\LANGpl{
\Question{
Znajdź zbiór rozwiązań następującej nierówności kwadratowej.
\[ 
 (x - 3)^{2}\geq 0
\]}
{\(\mathbb{R}\)}{\(\emptyset \)}{\(\mathbb{R}\setminus \{3\}\)}{\(\{3\}\)}{}{}}
\LANGcs{
\Question{
Určete množinu všech řešení dané nerovnice.
\[(x - 3)^{2}\geq 0\]}
{\(\mathbb{R}\)}{\(\emptyset \)}{\(\mathbb{R}\setminus \{3\}\)}{\(\{3\}\)}{}{}}
\LANGsk{
\Question{
Nájdite množinu všetkých riešení danej kvadratickej nerovnice. 
\[ 
 (x - 3)^{2}\geq 0
\]}
{\(\mathbb{R}\)}{\(\emptyset \)}{\(\mathbb{R}\setminus \{3\}\)}{\(\{3\}\)}{}{}}
\LANGes{
\Question{
Encuentra el conjunto de soluciones de la siguiente inecuación cuadrática.  
\[ 
 (x - 3)^{2}\geq 0
\]}
{\(\mathbb{R}\)}{\(\emptyset \)}{\(\mathbb{R}\setminus \{3\}\)}{\(\{3\}\)}{}{}}
\End

\Data\ID{9000039004}
\Updated{2021-08-19 03:08:16}
\Level{B}
\SubArea{Quadratic equations and inequalities}
\LANGen{
\Question{
Find all \(x\) such that the expression \( (x - 1)(x - 7) \) is negative.}
{\(x\in (1;7)\)}{\(x\in (-\infty ;1)\cup (7;+\infty )\)}{No such \(x\) 
exists.}{\(x\in \mathbb{R}\)}{}{}}
\LANGpl{
\Question{
Znajdź wszystkie wartości  \(x\), dla których podane wyrażenie jest ujemne.
\[ 
 (x - 1)(x - 7)
\]}
{\(x\in (1;7)\)}{\(x\in (-\infty ;1)\cup (7;+\infty )\)}{Takie \(x\) 
nie istnieje.}{\(x\in \mathbb{R}\)}{}{}}
\LANGcs{
\Question{
Pro která \(x\) 
nabývá součin \((x - 1)(x - 7)\) 
záporných hodnot?}
{\(x\in (1;7)\)}{\(x\in (-\infty ;1)\cup (7;+\infty )\)}{Takové \(x\) neexistuje.}{\(x\in \mathbb{R}\)}{}{}}
\LANGsk{
\Question{
Nájdite všetky hodnoty \(x\), pre ktoré je daný výraz záporný.
\[ 
 (x - 1)(x - 7)
\]}
{\(x\in (1;7)\)}{\(x\in (-\infty ;1)\cup (7;+\infty )\)}{Také \(x\) neexistuje.}{\(x\in \mathbb{R}\)}{}{}}
\LANGes{
\Question{
Halla todos los valores de \(x\) 
para los que la siguiente expresión es negativa.  
\[ 
 (x - 1)(x - 7)
\]}
{\(x\in (1;7)\)}{\(x\in (-\infty ;1)\cup (7;+\infty )\)}{Este valor de \(x\) 
no existe.}{\(x\in \mathbb{R}\)}{}{}}
\End

\Data\ID{1003067802}
\Updated{2020-08-05 05:08:43}
\Level{C}
\SubArea{Quadratic equations and inequalities}
\LANGen{
\Question{
For  \( x\in (-\infty;3) \)   choose the correct form of the equation 
\[ \left|x^2- 9 x + 20\right|=1 \]
that does not contain an absolute value.}
{\( x^2-9x+20=1 \)}{\( -x^2+9x-20=1 \)}{\( x^2-9x+20=-1 \)}{\( x^2+9x-20=1 \)}{}{}}
\LANGpl{
\Question{
Dla \( x\in (-\infty;3) \)   wybierz poprawny wzór równania
\[ \left|x^2- 9 x + 20\right|=1, \] który nie zawiera wartości bezwzględnej.}
{\( x^2-9x+20=1 \)}{\( -x^2+9x-20=1 \)}{\( x^2-9x+20=-1 \)}{\( x^2+9x-20=1 \)}{}{}}
\LANGcs{
\Question{
Vyberte správný tvar dané rovnice pro \( x\in (-\infty;3) \).
\[ \left|x^2- 9 x + 20\right|=1 \]}
{\( x^2-9x+20=1 \)}{\( -x^2+9x-20=1 \)}{\( x^2-9x+20=-1 \)}{\( x^2+9x-20=1 \)}{}{}}
\LANGsk{
\Question{
Pre  \( x\in (-\infty;3) \)   vyberte správny tvar danej rovnice. 
\[ \left|x^2- 9 x + 20\right|=1 \]}
{\( x^2-9x+20=1 \)}{\( -x^2+9x-20=1 \)}{\( x^2-9x+20=-1 \)}{\( x^2+9x-20=1 \)}{}{}}
\LANGes{
\Question{
Elige para \( x\in (-\infty;3) \) la forma correcta de la ecuación que no contenga un valor absoluto. 
\[ \left|x^2- 9 x + 20\right|=1 \].}
{\( x^2-9x+20=1 \)}{\( -x^2+9x-20=1 \)}{\( x^2-9x+20=-1 \)}{\( x^2+9x-20=1 \)}{}{}}
\End

\Data\ID{1003067803}
\Updated{2020-08-05 06:08:25}
\Level{C}
\SubArea{Quadratic equations and inequalities}
\LANGen{
\Question{
For \( x\in(1;2) \)   choose the correct form of the equation
\[ \left|2x^2-5x-7\right|=\left|x^2-3x+1\right| \]
that does not contain an absolute value.}
{\( -2x^2+5x+7=-x^2+3x-1 \)}{\( 2x^2-5x-7=-x^2+3x-1 \)}{\( 2x^2+5x+7=x^2-3x+1 \)}{\( -2x^2+5x+7=x^2-3x+1 \)}{}{}}
\LANGpl{
\Question{
Dla \( x\in(1;2) \)  wybierz poprawny wzór równania
\[ \left|2x^2-5x-7\right|=\left|x^2-3x+1\right|, \] który nie zawiera wartości bezwzględnej.}
{\( -2x^2+5x+7=-x^2+3x-1 \)}{\( 2x^2-5x-7=-x^2+3x-1 \)}{\( 2x^2+5x+7=x^2-3x+1 \)}{\( -2x^2+5x+7=x^2-3x+1 \)}{}{}}
\LANGcs{
\Question{
Vyberte správný tvar dané rovnice pro \( x\in(1;2) \).
\[ \left|2x^2-5x-7\right|=\left|x^2-3x+1\right| \]}
{\( -2x^2+5x+7=-x^2+3x-1 \)}{\( 2x^2-5x-7=-x^2+3x-1 \)}{\( 2x^2+5x+7=x^2-3x+1 \)}{\( -2x^2+5x+7=x^2-3x+1 \)}{}{}}
\LANGsk{
\Question{
Pre \( x\in(1;2) \)   vyberte správny tvar danej rovnice. 
\[ \left|2x^2-5x-7\right|=\left|x^2-3x+1\right| \]}
{\( -2x^2+5x+7=-x^2+3x-1 \)}{\( 2x^2-5x-7=-x^2+3x-1 \)}{\( 2x^2+5x+7=x^2-3x+1 \)}{\( -2x^2+5x+7=x^2-3x+1 \)}{}{}}
\LANGes{
\Question{
Elige para \( x\in(1;2) \) la forma correcta de la ecuación que no contenga un valor absoluto
\[ \left|2x^2-5x-7\right|=\left|x^2-3x+1\right| \].}
{\( -2x^2+5x+7=-x^2+3x-1 \)}{\( 2x^2-5x-7=-x^2+3x-1 \)}{\( 2x^2+5x+7=x^2-3x+1 \)}{\( -2x^2+5x+7=x^2-3x+1 \)}{}{}}
\End

\Data\ID{1003067804}
\Updated{2020-08-05 05:08:23}
\Level{C}
\SubArea{Quadratic equations and inequalities}
\LANGen{
\Question{
For \( x\in[4;\infty) \) choose the correct form of the equation 
\[ \left|-x^2+3x+4\right|=\left|-2 x^2+ 11 x - 12\right| \]
that does not contain an absolute value.}
{\( x^2-3x-4=2x^2-11x+12 \)}{\( x^2-3x-4=-2x^2+11x-12 \)}{\(-x^2+3x+4=2x^2-11x+12 \)}{\( -x^2+3x+4=-2x^2+11x-12 \)}{}{}}
\LANGpl{
\Question{
Dla\( x\in\langle4;\infty) \) wybierz poprawny wzór równania
\[ \left|-x^2+3x+4\right|=\left|-2 x^2+ 11 x - 12\right|, \]  który nie zawiera wartości bezwzględnej.}
{\( x^2-3x-4=2x^2-11x+12 \)}{\( x^2-3x-4=-2x^2+11x-12 \)}{\(-x^2+3x+4=2x^2-11x+12 \)}{\( -x^2+3x+4=-2x^2+11x-12 \)}{}{}}
\LANGcs{
\Question{
Vyberte správný tvar dané rovnice pro \( x\in\langle4;\infty) \).
\[ \left|-x^2+3x+4\right|=\left|-2 x^2+ 11 x - 12\right| \]}
{\( x^2-3x-4=2x^2-11x+12 \)}{\( x^2-3x-4=-2x^2+11x-12 \)}{\(-x^2+3x+4=2x^2-11x+12 \)}{\( -x^2+3x+4=-2x^2+11x-12 \)}{}{}}
\LANGsk{
\Question{
Pre \( x\in\langle4;\infty) \) vyberte správny tvar danej rovnice. 
\[ \left|-x^2+3x+4\right|=\left|-2 x^2+ 11 x - 12\right| \]}
{\( x^2-3x-4=2x^2-11x+12 \)}{\( x^2-3x-4=-2x^2+11x-12 \)}{\(-x^2+3x+4=2x^2-11x+12 \)}{\( -x^2+3x+4=-2x^2+11x-12 \)}{}{}}
\LANGes{
\Question{
Elige para \( x\in[4;\infty) \) la forma correcta de la ecuación que no contenga un valor absoluto  
\[ \left|-x^2+3x+4\right|=\left|-2 x^2+ 11 x - 12\right| \].}
{\( x^2-3x-4=2x^2-11x+12 \)}{\( x^2-3x-4=-2x^2+11x-12 \)}{\(-x^2+3x+4=2x^2-11x+12 \)}{\( -x^2+3x+4=-2x^2+11x-12 \)}{}{}}
\End

\Data\ID{1003067805}
\Updated{2020-08-05 05:08:41}
\Level{C}
\SubArea{Quadratic equations and inequalities}
\LANGen{
\Question{
For  \( x\in[-3;5] \)  find the solution set of the following equation.
\[ \left|(x+3)(x-5)\right|=5 \]}
{\( \left\{ 1-\sqrt{11};1+\sqrt{11} \right\} \)}{\( \left\{ 1-\sqrt{21};1+\sqrt{21} \right\} \)}{\( \{ -3; 5 \} \)}{\( \left\{1-\sqrt{21}; 1-\sqrt{11};1+\sqrt{11};1+\sqrt{21} \right\} \)}{}{}}
\LANGpl{
\Question{
Dla \( x\in\langle-3;5\rangle \)  wyznacz zbiór rozwiązań następującego równania.
\[ \left|(x+3)(x-5)\right|=5 \]}
{\( \left\{ 1-\sqrt{11};1+\sqrt{11} \right\} \)}{\( \left\{ 1-\sqrt{21};1+\sqrt{21} \right\} \)}{\( \{ -3; 5 \} \)}{\( \left\{1-\sqrt{21}; 1-\sqrt{11};1+\sqrt{11};1+\sqrt{21} \right\} \)}{}{}}
\LANGcs{
\Question{
Určete množinu řešení dané rovnice pro  \( x\in\langle-3;5\rangle \).
\[ \left|(x+3)(x-5)\right|=5 \]}
{\( \left\{ 1-\sqrt{11};1+\sqrt{11} \right\} \)}{\( \left\{ 1-\sqrt{21};1+\sqrt{21} \right\} \)}{\( \{ -3; 5 \} \)}{\( \left\{1-\sqrt{21}; 1-\sqrt{11};1+\sqrt{11};1+\sqrt{21} \right\} \)}{}{}}
\LANGsk{
\Question{
Pre  \( x\in\langle-3;5\rangle \)  určte množinu riešení danej rovnice.
\[ \left|(x+3)(x-5)\right|=5 \]}
{\( \left\{ 1-\sqrt{11};1+\sqrt{11} \right\} \)}{\( \left\{ 1-\sqrt{21};1+\sqrt{21} \right\} \)}{\( \{ -3; 5 \} \)}{\( \left\{1-\sqrt{21}; 1-\sqrt{11};1+\sqrt{11};1+\sqrt{21} \right\} \)}{}{}}
\LANGes{
\Question{
Encuentra el conjunto de soluciones de la siguiente ecuación para \( x\in[-3;5] \). 
\[ \left|(x+3)(x-5)\right|=5 \]}
{\( \left\{ 1-\sqrt{11};1+\sqrt{11} \right\} \)}{\( \left\{ 1-\sqrt{21};1+\sqrt{21} \right\} \)}{\( \{ -3; 5 \} \)}{\( \left\{1-\sqrt{21}; 1-\sqrt{11};1+\sqrt{11};1+\sqrt{21} \right\} \)}{}{}}
\End

\Data\ID{1003067806}
\Updated{2020-08-09 02:08:01}
\Level{C}
\SubArea{Quadratic equations and inequalities}
\LANGen{
\Question{
Find the solution set of the following equation.
\[ |(x-1)(x+2)|=4 \]}
{\( \{ -3; 2 \} \)}{\( \{ 2 \} \)}{\( \{ -2; 3 \} \)}{\( \{ -2 \} \)}{}{}}
\LANGpl{
\Question{
Wyznacz zbiór rozwiązań następującego równania. 
\[ |(x-1)(x+2)|=4 \]}
{\( \{ -3; 2 \} \)}{\( \{ 2 \} \)}{\( \{ -2; 3 \} \)}{\( \{ -2 \} \)}{}{}}
\LANGcs{
\Question{
Určete množinu řešení dané rovnice.
\[ |(x-1)(x+2)|=4 \]}
{\( \{ -3; 2 \} \)}{\( \{ 2 \} \)}{\( \{ -2; 3 \} \)}{\( \{ -2 \} \)}{}{}}
\LANGsk{
\Question{
Určte množinu riešení danej rovnice.
\[ |(x-1)(x+2)|=4 \]}
{\( \{ -3; 2 \} \)}{\( \{ 2 \} \)}{\( \{ -2; 3 \} \)}{\( \{ -2 \} \)}{}{}}
\LANGes{
\Question{
Encuentra el conjunto de soluciones de la siguiente ecuación. 
\[ |(x-1)(x+2)|=4 \]}
{\( \{ -3; 2 \} \)}{\( \{ 2 \} \)}{\( \{ -2; 3 \} \)}{\( \{ -2 \} \)}{}{}}
\End

\Data\ID{1003067807}
\Updated{2020-08-09 03:08:00}
\Level{C}
\SubArea{Quadratic equations and inequalities}
\LANGen{
\Question{
Find the solution set of the following equation.
\[ 2x^2+ 4x - 30=\left|2 x^2+ 4 x - 30\right| \]}
{\( (-\infty;-5]\cup[3;\infty) \)}{\( [-5;3] \)}{\( (-\infty;-3]\cup[5;\infty) \)}{\( [-3;5] \)}{}{}}
\LANGpl{
\Question{
Wyznacz zbiór rozwiązań następującego równania. 
\[ 2x^2+ 4x - 30=\left|2 x^2+ 4 x - 30\right| \]}
{\( (-\infty;-5\rangle\cup\langle3;\infty) \)}{\( \langle-5;3\rangle \)}{\( (-\infty;-3\rangle\cup\langle5;\infty) \)}{\( \langle-3;5\rangle \)}{}{}}
\LANGcs{
\Question{
Určete množinu řešení dané rovnice.
\[ 2x^2+ 4x - 30=\left|2 x^2+ 4 x - 30\right| \]}
{\( (-\infty;-5\rangle\cup\langle3;\infty) \)}{\( \langle-5;3\rangle \)}{\( (-\infty;-3\rangle\cup\langle5;\infty) \)}{\( \langle-3;5\rangle \)}{}{}}
\LANGsk{
\Question{
Určte množinu riešení danej rovnice.
\[ 2x^2+ 4x - 30=\left|2 x^2+ 4 x - 30\right| \]}
{\( (-\infty;-5\rangle\cup\langle3;\infty) \)}{\( \langle-5;3\rangle \)}{\( (-\infty;-3\rangle\cup\langle5;\infty) \)}{\( \langle-3;5\rangle \)}{}{}}
\LANGes{
\Question{
Encuentra el conjunto de soluciones de la siguiente ecuación. 
\[ 2x^2+ 4x - 30=\left|2 x^2+ 4 x - 30\right| \]}
{\( (-\infty;-5]\cup[3;\infty) \)}{\( [-5;3] \)}{\( (-\infty;-3]\cup[5;\infty) \)}{\( [-3;5] \)}{}{}}
\End

\Data\ID{1003067808}
\Updated{2020-08-09 03:08:30}
\Level{C}
\SubArea{Quadratic equations and inequalities}
\LANGen{
\Question{
Find the solution set of the following equation.
\[ -2 x^2+ 5 x + 3 =\left|-2 x^2+ 5 x + 3\right| \]}
{\( \left[-\frac12;3\right] \)}{\( [-5;3]  \)}{\( \left(-\infty;-\frac12\right]\cup[3;\infty) \)}{\( (-\infty,-3]\cup[5,\infty) \)}{}{}}
\LANGpl{
\Question{
Wyznacz zbiór rozwiązań następującego równania. 
\[ -2 x^2+ 5 x + 3 =\left|-2 x^2+ 5 x + 3\right| \]}
{\( \left\langle-\frac12;3\right\rangle \)}{\( \langle-5;3\rangle  \)}{\( \left(-\infty;-\frac12\right\rangle\cup\langle3;\infty) \)}{\( (-\infty,-3\rangle\cup\langle5,\infty) \)}{}{}}
\LANGcs{
\Question{
Určete množinu řešení dané rovnice.
\[ -2 x^2+ 5 x + 3 =\left|-2 x^2+ 5 x + 3\right| \]}
{\( \left\langle-\frac12;3\right\rangle \)}{\( \langle-5;3\rangle  \)}{\( \left(-\infty;-\frac12\right\rangle\cup\langle3;\infty) \)}{\( (-\infty,-3\rangle\cup\langle5,\infty) \)}{}{}}
\LANGsk{
\Question{
Určte množinu riešení danej rovnice.
\[ -2 x^2+ 5 x + 3 =\left|-2 x^2+ 5 x + 3\right| \]}
{\( \left\langle-\frac12;3\right\rangle \)}{\( \langle-5;3\rangle  \)}{\( \left(-\infty;-\frac12\right\rangle\cup\langle3;\infty) \)}{\( (-\infty,-3\rangle\cup\langle5,\infty) \)}{}{}}
\LANGes{
\Question{
Encuentra el conjunto de soluciones de la siguiente ecuación. 
\[ -2 x^2+ 5 x + 3 =\left|-2 x^2+ 5 x + 3\right| \]}
{\( \left[-\frac12;3\right] \)}{\( [-5;3]  \)}{\( \left(-\infty;-\frac12\right]\cup[3;\infty) \)}{\( (-\infty,-3]\cup[5,\infty) \)}{}{}}
\End

\Data\ID{1003067810}
\Updated{2020-08-09 03:08:05}
\Level{C}
\SubArea{Quadratic equations and inequalities}
\LANGen{
\Question{
Find the solution set of the following equation.
\[ |x-4|\cdot(x+4)=4 \]}
{\( \left\{-2\sqrt3;2\sqrt3;2\sqrt5\right\} \)}{\( \{ -4; 4 \} \)}{\( \left\{ -2\sqrt3;2\sqrt3\right\} \)}{\( \left\{ 2\sqrt3;2\sqrt5\right\} \)}{\( \left\{-2\sqrt5;-2\sqrt3;2\sqrt3;2\sqrt5 \right\} \)}{}}
\LANGpl{
\Question{
Wyznacz zbiór rozwiązań następującego równania. 
\[ |x-4|\cdot(x+4)=4 \]}
{\( \left\{-2\sqrt3;2\sqrt3;2\sqrt5\right\} \)}{\( \{ -4; 4 \} \)}{\( \left\{ -2\sqrt3;2\sqrt3\right\} \)}{\( \left\{ 2\sqrt3;2\sqrt5\right\} \)}{\( \left\{-2\sqrt5;-2\sqrt3;2\sqrt3;2\sqrt5 \right\} \)}{}}
\LANGcs{
\Question{
Určete množinu řešení dané rovnice.
\[ |x-4|\cdot(x+4)=4 \]}
{\( \left\{-2\sqrt3;2\sqrt3;2\sqrt5\right\} \)}{\( \{ -4; 4 \} \)}{\( \left\{ -2\sqrt3;2\sqrt3\right\} \)}{\( \left\{ 2\sqrt3;2\sqrt5\right\} \)}{\( \left\{-2\sqrt5;-2\sqrt3;2\sqrt3;2\sqrt5 \right\} \)}{}}
\LANGsk{
\Question{
Určte množinu riešení danej rovnice.
\[ |x-4|\cdot(x+4)=4 \]}
{\( \left\{-2\sqrt3;2\sqrt3;2\sqrt5\right\} \)}{\( \{ -4; 4 \} \)}{\( \left\{ -2\sqrt3;2\sqrt3\right\} \)}{\( \left\{ 2\sqrt3;2\sqrt5\right\} \)}{\( \left\{-2\sqrt5;-2\sqrt3;2\sqrt3;2\sqrt5 \right\} \)}{}}
\LANGes{
\Question{
Encuentra el conjunto de soluciones de la siguiente ecuación. 
\[ |x-4|\cdot(x+4)=4 \]}
{\( \left\{-2\sqrt3;2\sqrt3;2\sqrt5\right\} \)}{\( \{ -4; 4 \} \)}{\( \left\{ -2\sqrt3;2\sqrt3\right\} \)}{\( \left\{ 2\sqrt3;2\sqrt5\right\} \)}{\( \left\{-2\sqrt5;-2\sqrt3;2\sqrt3;2\sqrt5 \right\} \)}{}}
\End

\Data\ID{1003085404}
\Updated{2020-08-20 07:08:36}
\Level{C}
\SubArea{Quadratic equations and inequalities}
\LANGen{
\Question{
The surface area of a cuboid is \( 11\,776\,\mathrm{cm}^2 \). Its dimensions are in the ratio \( 3:5:1 \). Find the volume of the cuboid.}
{\( 61\,440\,\mathrm{cm}^3 \)}{\( 15\,360\,\mathrm{cm}^3 \)}{\( 30\,720\,\mathrm{cm}^3 \)}{\( 21\,722\,\mathrm{cm}^3 \)}{}{}}
\LANGpl{
\Question{
Pole powierzchni prostopadłościanu wynosi \( 11\,776\,\mathrm{cm}^2 \). Jego wymiary podano w stosunku \( 3:5:1 \). Oblicz objętość tego prostopadłościanu.}
{\( 61\,440\,\mathrm{cm}^3 \)}{\( 15\,360\,\mathrm{cm}^3 \)}{\( 30\,720\,\mathrm{cm}^3 \)}{\( 21\,722\,\mathrm{cm}^3 \)}{}{}}
\LANGcs{
\Question{
Povrch kvádru je \( 11\,776\,\mathrm{cm}^2 \). Délky jeho stran jsou v poměru \( 3:5:1 \). Určete objem daného kvádru.}
{\( 61\,440\,\mathrm{cm}^3 \)}{\( 15\,360\,\mathrm{cm}^3 \)}{\( 30\,720\,\mathrm{cm}^3 \)}{\( 21\,722\,\mathrm{cm}^3 \)}{}{}}
\LANGsk{
\Question{
Povrch kvádra je \( 11\,776\,\mathrm{cm}^2 \). Jeho rozmery sú v pomere \( 3:5:1 \). Určte objem daného kvádra.}
{\( 61\,440\,\mathrm{cm}^3 \)}{\( 15\,360\,\mathrm{cm}^3 \)}{\( 30\,720\,\mathrm{cm}^3 \)}{\( 21\,722\,\mathrm{cm}^3 \)}{}{}}
\LANGes{
\Question{
El área de un ortoedro es \( 11\,776\,\mathrm{cm}^2 \). Sus lados están en la proporción \( 3:5:1 \). Encuentra el volumen de este ortoedro.}
{\( 61\,440\,\mathrm{cm}^3 \)}{\( 15\,360\,\mathrm{cm}^3 \)}{\( 30\,720\,\mathrm{cm}^3 \)}{\( 21\,722\,\mathrm{cm}^3 \)}{}{}}
\End

\Data\ID{1003085405}
\Updated{2020-08-22 04:08:39}
\Level{C}
\SubArea{Quadratic equations and inequalities}
\LANGen{
\Question{
Little Red Riding Hood ran (at instantaneous speed) through the forest to see her grandmother, who lives in a cottage, which is \( 4\,\mathrm{km} \) distant. If she ran by \( 4\,\mathrm{km/h} \) faster, she would meet her grandmother \( 10 \) minutes sooner. What was Little Red Riding Hood’s speed?}
{\( 8\,\mathrm{km/h} \)}{\( 12\,\mathrm{km/h} \)}{\( 10\,\mathrm{km/h} \)}{\( 6\,\mathrm{km/h} \)}{}{}}
\LANGpl{
\Question{
Czerwony Kapturek biegł (z chwilową prędkością) przez las do swojej babci, która mieszka w chatce oddalonej o \( 4\,\mathrm{km} \).  Gdyby biegła o  \( 4\,\mathrm{km/h} \) szybciej, spotkałaby się z babcią  \( 10 \) wcześniej. Z jaką prędkością biegł Czerwony Kapturek?}
{\( 8\,\mathrm{km/h} \)}{\( 12\,\mathrm{km/h} \)}{\( 10\,\mathrm{km/h} \)}{\( 6\,\mathrm{km/h} \)}{}{}}
\LANGcs{
\Question{
Karkulka běžela (stálou rychlostí) přes les za babičkou, která bydlí v chalupě vzdálené \( 4\,\mathrm{km} \). Kdyby běžela o \( 4\,\mathrm{km/h} \) rychleji, cesta by jí trvala o  \( 10 \) minut méně. Jakou rychlostí Karkulka běžela?}
{\( 8\,\mathrm{km/h} \)}{\( 12\,\mathrm{km/h} \)}{\( 10\,\mathrm{km/h} \)}{\( 6\,\mathrm{km/h} \)}{}{}}
\LANGsk{
\Question{
Červená Čiapočka bežala (stálou rýchlosťou) cez les navštíviť babičku, ktorá býva v chalupe vzdialenej \( 4\,\mathrm{km} \). Keby bežala o \( 4\,\mathrm{km/h} \) rýchlejšie, stretla by sa so svojou babičkou o \( 10 \) minút skôr. Akou rýchlosťou Červená Čiapočka bežala?}
{\( 8\,\mathrm{km/h} \)}{\( 12\,\mathrm{km/h} \)}{\( 10\,\mathrm{km/h} \)}{\( 6\,\mathrm{km/h} \)}{}{}}
\LANGes{
\Question{
Caperucita Roja corrió (a una velocidad continua) por el bosque para visitar a su abuela que vivía en su casita a \( 4\,\mathrm{km} \) de la casa de Caperucita. Si corriera \( 4\,\mathrm{km/h} \) más rápido, llegaría a la casita de su abuela \( 10 \) minutos antes. ¿Qué velocidad llevó Caperucita?}
{\( 8\,\mathrm{km/h} \)}{\( 12\,\mathrm{km/h} \)}{\( 10\,\mathrm{km/h} \)}{\( 6\,\mathrm{km/h} \)}{}{}}
\End

\Data\ID{1003085408}
\Updated{2020-08-22 04:08:53}
\Level{C}
\SubArea{Quadratic equations and inequalities}
\LANGen{
\Question{
A swimming pool can be filled by two pipes in \( 5 \) hours. It takes \( 24 \) hours longer to fill the pool by only the first pipe than by only the second one. In how many hours does the first pipe fill the pool, in how many hours does the second one? Find the sum of both times.}
{\( 36 \) hours}{\( 20 \) hours}{\( 18 \) hours}{\( 32 \) hours}{}{}}
\LANGpl{
\Question{
Dwie rury są w stanie napełnić basen w ciągu \( 5 \) godzin. Pierwsza rura napełnia basen o  \( 24 \) godziny  dłużej niż druga. Ile godzin zajmie każdej z tych rur z osobna napełnienie tego  basenu. Oblicz sumę obydwu czasów.}
{\( 36 \) godzin}{\( 20 \) godzin}{\( 18 \) godzin}{\( 32 \) godzin}{}{}}
\LANGcs{
\Question{
Bazén se dvěma přítoky otevřenými současně napustí za \( 5 \) hodin. Pokud by byl otevřen pouze první přítok, napouštěl by se bazén o  \( 24 \) hodin déle než v případě, že by byl otevřen pouze druhý přítok. Za jak dlouho se napustí bazén v případě, že je otevřen pouze první přítok a za jak dlouho, je-li otevřen pouze druhý přítok? Vyberte odpověď odpovídající součtu obou časů.}
{\( 36 \) hodin}{\( 20 \) hodin}{\( 18 \) hodin}{\( 32 \) hodin}{}{}}
\LANGsk{
\Question{
Bazén sa napustí dvoma prítokmi otvorenými súčasne za \( 5 \) hodín. Ak by bol otvorený len prvý prítok, napúšťal by sa bazén o \( 24 \) hodín dlhšie než v prípade, ak by bol otvorený len druhý prítok. Za ako dlho sa napustí bazén v prípade, že je otvorený len prvý prítok a za ako dlho, ak je otvorený len druhý prítok?   Vyberte odpoveď zodpovedajúcu súčtu obidvoch časov.}
{\( 36 \) hodín}{\( 20 \) hodín}{\( 18 \) hodín}{\( 32 \) hodín}{}{}}
\LANGes{
\Question{
Una piscina puede estar llena en \( 5 \) horas usando dos chorros. Utilizando solo el primer chorro llenar la piscina tarda \( 24 \) horas más que utilizar solo el segundo. ¿Cuánto tiempo tarda llenar la piscina con el primer chorro y cuánto tiempo tarda llenarla con el segundo? Calcula la suma de estos dos tiempos.}
{\( 36 \) horas}{\( 20 \) horas}{\( 18 \) horas}{\( 32 \) horas}{}{}}
\End

\Data\ID{1003162910}
\Updated{2020-08-09 03:08:13}
\Level{C}
\SubArea{Quadratic equations and inequalities}
\LANGen{
\Question{
Find the sum of all the roots of the equation \( \left|x^2+2x\right|-6=|3x| \).}
{\( -3 \)}{\( 9 \)}{\( 0 \)}{\( -2 \)}{}{}}
\LANGpl{
\Question{
Wyznacz sumę wszystkich pierwiastków równania \( \left|x^2+2x\right|-6=|3x| \).}
{\( -3 \)}{\( 9 \)}{\( 0 \)}{\( -2 \)}{}{}}
\LANGcs{
\Question{
Určete součet všech kořenů rovnice \( \left|x^2+2x\right|-6=|3x| \).}
{\( -3 \)}{\( 9 \)}{\( 0 \)}{\( -2 \)}{}{}}
\LANGsk{
\Question{
Určte súčet všetkých koreňov rovnice \( \left|x^2+2x\right|-6=|3x| \).}
{\( -3 \)}{\( 9 \)}{\( 0 \)}{\( -2 \)}{}{}}
\LANGes{
\Question{
Encuentra la suma de las raíces de la ecuación \( \left|x^2+2x\right|-6=|3x| \).}
{\( -3 \)}{\( 9 \)}{\( 0 \)}{\( -2 \)}{}{}}
\End

\Data\ID{1003187313}
\Updated{2020-08-20 11:08:34}
\Level{C}
\SubArea{Quadratic equations and inequalities}
\LANGen{
\Question{
Which of the given sets contains all negative integers that satisfy the inequality  \( \sqrt{(2x-8)^2} < 14 \)?}
{\( \{-2;-1\} \)}{\( \{-3;-2;-1\} \)}{\( \{-10;-9;-8;-7;-6;-5;-4;-3;-2;-1\} \)}{\( \{-4;-3;-2;-1\} \)}{}{}}
\LANGpl{
\Question{
Który z podanych zbiorów zawiera wszystkie ujemne liczby całkowite, które spełniają nierówność  \( \sqrt{(2x-8)^2} < 14 \)?}
{\( \{-2;-1\} \)}{\( \{-3;-2;-1\} \)}{\( \{-10;-9;-8;-7;-6;-5;-4;-3;-2;-1\} \)}{\( \{-4;-3;-2;-1\} \)}{}{}}
\LANGcs{
\Question{
Která z daných množin obsahuje všechna celá záporná řešení dané nerovnice?  \[ \sqrt{(2x-8)^2} < 14 \]}
{\( \{-2;-1\} \)}{\( \{-3;-2;-1\} \)}{\( \{-10;-9;-8;-7;-6;-5;-4;-3;-2;-1\} \)}{\( \{-4;-3;-2;-1\} \)}{}{}}
\LANGsk{
\Question{
Ktorá z daných množín obsahuje všetký celé záporné riešenia dané nerovnice? \[ \sqrt{(2x-8)^2} < 14 \]}
{\( \{-2;-1\} \)}{\( \{-3;-2;-1\} \)}{\( \{-10;-9;-8;-7;-6;-5;-4;-3;-2;-1\} \)}{\( \{-4;-3;-2;-1\} \)}{}{}}
\LANGes{
\Question{
¿Cuál de los siguientes conjuntos contiene todos los números enteros negativos que son soluciones de la siguiente inecuación \( \sqrt{(2x-8)^2} < 14 \)?}
{\( \{-2;-1\} \)}{\( \{-3;-2;-1\} \)}{\( \{-10;-9;-8;-7;-6;-5;-4;-3;-2;-1\} \)}{\( \{-4;-3;-2;-1\} \)}{}{}}
\End

\Data\ID{9000033708}
\Updated{2020-08-22 05:08:35}
\Level{C}
\SubArea{Quadratic equations and inequalities}
\LANGen{
\Question{
A stone has been thrown vertically up at the velocity
\(15\, \mathrm{m}\, \mathrm{s}^{-1}\) from the
initial height \(10\, \mathrm{m}\).
How long (in seconds) has been the height of the stone at least
\(20\, \mathrm{m}\)? 

Hint: The height \(h\) is
given by the expression \(h = s_{0} + v_{0}t -\frac{1} 
{2}gt^{2}\), the
standard acceleration is \(g\mathop{\mathop{\doteq }}\nolimits 10\, \mathrm{m}\, \mathrm{s}^{-2}\).}
{exactly \(1\, \mathrm{s}\)}{less than \(1\, \mathrm{s}\)}{more than \(1\, \mathrm{s}\)}{The information is not sufficient to give a definite answer.}{}{}}
\LANGpl{
\Question{
Kamień został rzucony pionowo w górę z prędkością
\(15\, \mathrm{m}\, \mathrm{s}^{-1}\) od początkowej wysokości \(10\, \mathrm{m}\).
Jak długo (w sekundach) wysokość kamienia miała przynajmniej  \(20\, \mathrm{m}\)?

Wskazówka: Wysokość  \(h\) jest
równa wyrażeniu  \(h = s_{0} + v_{0}t -\frac{1} 
{2}gt^{2}\), standardowe przyspieszenie wynosi  \(g\mathop{\mathop{\doteq }}\nolimits 10\, \mathrm{m}\, \mathrm{s}^{-2}\).}
{dokładnie \(1\, \mathrm{s}\)}{mniej niż \(1\, \mathrm{s}\)}{więcej niż \(1\, \mathrm{s}\)}{Informacje nie są wystarczające, by udzielić jednoznacznej  odpowiedzi.}{}{}}
\LANGcs{
\Question{
Kámen byl ve výšce \(10\, \mathrm{m}\) 
nad zemí vržen svisle vzhůru rychlostí
\(15\, \mathrm{m}\, \mathrm{s}^{-1}\).
Rozhodněte, jak dlouho byla jeho poloha ve výšce alespoň
\(20\, \mathrm{m}\) nad
zemí? 

Nápověda: Pro výšku \(h\) 
využijte vztah \(h = s_{0} + v_{0}t -\frac{1} 
{2}gt^{2}\), za hodnotu
tíhového zrychlení dosaďte \(g\mathop{\mathop{\doteq }}\nolimits 10\, \mathrm{m}\, \mathrm{s}^{-2}\).}
{právě \(1\, \mathrm{s}\)}{méně než \(1\, \mathrm{s}\)}{déle než \(1\, \mathrm{s}\)}{Ze zadání nelze určit.}{}{}}
\LANGsk{
\Question{
Kameň bol  vrhnutý zvislo nahor rýchlosťou 
\(15\, \mathrm{m}\, \mathrm{s}^{-1}\) vo výške
 \(10\, \mathrm{m}\) nad zemou. Rozhodnite, ako dlho (v sekundách) bola jeho poloha vo výške aspoň 
\(20\, \mathrm{m}\) nad zemou.

Pomôcka: Pre výšku  \(h\) využite vzťah
 \(h = s_{0} + v_{0}t -\frac{1} 
{2}gt^{2}\), hodnota gravitačného zrýchlenia je \(g\mathop{\mathop{\doteq }}\nolimits 10\, \mathrm{m}\, \mathrm{s}^{-2}\).}
{presne \(1\, \mathrm{s}\)}{menej ako \(1\, \mathrm{s}\)}{viac ako \(1\, \mathrm{s}\)}{Informácie nie sú dostatočné na poskytnutie jednoznačné odpovede.}{}{}}
\LANGes{
\Question{
Una piedra fue tirada verticalmente hacia arriba desde una altura de \(10\, \mathrm{m}\) con una velocidad de \(15\, \mathrm{m}\, \mathrm{s}^{-1}\) .
¿Cuánto tiempo (en segundos) estuvo a una altura de mínima de \(20\, \mathrm{m}\)? 

Ayuda: La altura \(h\) se expresa \(h = s_{0} + v_{0}t -\frac{1} 
{2}gt^{2}\), la gravedad de la Tierra es \(g\mathop{\mathop{\doteq }}\nolimits 10\, \mathrm{m}\, \mathrm{s}^{-2}\).}
{exactamente \(1\, \mathrm{s}\)}{menos de \(1\, \mathrm{s}\)}{más de \(1\, \mathrm{s}\)}{No hay suficientes datos para poder responder.}{}{}}
\End

\Data\ID{9000033709}
\Updated{2020-08-22 05:08:35}
\Level{C}
\SubArea{Quadratic equations and inequalities}
\LANGen{
\Question{
A square shaped garden with the side
\(a\) should be reduced
by a length \(x\) to
another square garden. The difference between the areas of the gardens should not be bigger than
\(25\%\) of the original area. Find
the possible values of \(x\).}
{\(x\leq a -\frac{\sqrt{3}} 
 {2} a\)}{\(x\leq \sqrt{3}a\)}{\(x\leq \frac{3}
{4}a\)}{\(x\leq a + \frac{\sqrt{3}} 
 {2} a\)}{}{}}
\LANGpl{
\Question{
Ogród w kształcie kwadratu o boku
\(a\) powinien zostać zmniejszony
o długość  \(x\) tak, by powstał inny kwadratowy ogród. 
 Różnica pomiędzy powierzchniami ogrodów nie powinna być większa niż
\(25\%\) początkowej powierzchni. Znajdź możliwe wartości \(x\).}
{\(x\leq a -\frac{\sqrt{3}} 
 {2} a\)}{\(x\leq \sqrt{3}a\)}{\(x\leq \frac{3}
{4}a\)}{\(x\leq a + \frac{\sqrt{3}} 
 {2} a\)}{}{}}
\LANGcs{
\Question{
Rozměry čtvercové parcely o délce strany
\(a\) je třeba
zmenšit o délku \(x\) 
tak, aby zůstal zachován její čtvercový půdorys a aby se její obsah
nezmenšil o více než jednu čtvrtinu původního obsahu. O jakou délku tedy
můžeme rozměr parcely zmenšit?}
{\(x\leq a -\frac{\sqrt{3}} 
 {2} a\)}{\(x\leq \sqrt{3}a\)}{\(x\leq \frac{3}
{4}a\)}{\(x\leq a + \frac{\sqrt{3}} 
 {2} a\)}{}{}}
\LANGsk{
\Question{
Rozmery štvorcovej záhrady s dĺžkou strany
\(a\) je treba zmenšiť o dĺžku
 \(x\) tak, aby zostal zachovaný jej štvorcový pôdorys a aby sa jej obsah nezmenšil o viac než 
\(25\%\) pôvodného obsahu.
O akú dĺžku \(x\) môžeme rozmer záhrady zmenšiť?}
{\(x\leq a -\frac{\sqrt{3}} 
 {2} a\)}{\(x\leq \sqrt{3}a\)}{\(x\leq \frac{3}
{4}a\)}{\(x\leq a + \frac{\sqrt{3}} 
 {2} a\)}{}{}}
\LANGes{
\Question{
Un jardín cuadrado cuyo lado es \(a\) debe ser reducido a \(x\) para formar otro jardín cuadrado. La diferencia entre las áreas de los jardines no puede ser más de un \(25\%\) del jardín original. Halla los valores posibles de \(x\).}
{\(x\leq a -\frac{\sqrt{3}} 
 {2} a\)}{\(x\leq \sqrt{3}a\)}{\(x\leq \frac{3}
{4}a\)}{\(x\leq a + \frac{\sqrt{3}} 
 {2} a\)}{}{}}
\End
